\documentclass[UTF8,12pt,a4paper,fontset=fandol]{ctexbook}

% 从项目根目录或本文件所在目录均可定位共享样式。
\makeatletter
\def\input@path{{../}}
\makeatother
\usepackage{researchnotes}

\title{数值分析与科学计算}
\author{}
\date{}

\begin{document}
\maketitle
\tableofcontents

\chapter{绪论：从数学问题到可靠计算}
\label{chap:na-intro}

\section{研究对象与基本问题}
\keyterm{数值分析（numerical analysis）}研究数学问题的近似计算方法及其误差、稳定性和计算代价。\keyterm{科学计算（scientific computing）}进一步把这些方法组织成能够解决科学与工程问题的计算过程，涉及模型、离散化、算法实现和结果验证。一个方程有唯一解，不等于有限精度计算能够准确地求出它；一个程序返回许多位小数，也不等于这些小数具有相应的有效精度。

以求解导热问题为例，物理规律首先给出温度满足的微分方程，边界条件和材料参数决定具体模型；空间网格把连续函数转化为有限维向量，时间步进把连续演化转化为递推，线性代数算法再求出每一步的近似值。测量误差、模型误差、离散误差、迭代误差和舍入误差发生在不同层次，应分别识别。减小网格宽度通常降低离散误差，却也增加计算代价，并可能使代数系统更难求解。

本书围绕三个相互联系的问题展开：所求数学对象是否存在、唯一并且对数据连续依赖；离散或迭代近似是否趋于该对象以及趋近得多快；有限精度实现是否保留了理论算法的可靠性。第一类问题属于问题本身的性质，后两类还取决于算法。选择算法时应同时考察精度要求、数据规模和可利用的结构，而不能只比较公式的阶数。

\begin{example}[很小的残差能说明什么]
\label{ex:na-intro-residual}
设 $0<\varepsilon\ll1$，考虑 $\varepsilon x=\varepsilon$，其精确解为 $x=1$。近似解 $\widehat x=0$ 的残差为 $\varepsilon$，绝对值很小，但解的绝对误差等于 $1$。这说明绝对残差必须结合问题尺度解释。若把等式两端同除以 $\varepsilon$，残差变成 $1$，而数学问题没有改变。后文将用归一化残差和条件数把这一直觉写成可核查的不等式。
\end{example}

\section{本科定位、知识依赖与阅读路线}
本书面向已学习一元与多元微积分、线性代数的本科生。常微分方程初值问题的基本概念在相应章节回顾；蒙特卡洛方法还需要期望、方差、独立性和中心极限定理。编程部分仅要求理解变量、循环、函数和数组。线性空间的完备性、矩阵谱分解等结果按需要说明，不要求预先学习泛函分析。

第~\ref{chap:na-errors} 章建立误差语言，第~\ref{chap:na-roots} 章讨论非线性方程；第~\ref{chap:na-direct} 至第~\ref{chap:na-eigen} 章组成数值线性代数单元；第~\ref{chap:na-interpolation}、\ref{chap:na-quadrature} 章研究函数的离散表示与数值微积分；第~\ref{chap:na-ode}、\ref{chap:na-pde} 章把这些方法用于微分方程；第~\ref{chap:na-fourier}、\ref{chap:na-montecarlo} 章介绍傅里叶计算和随机积分；第~\ref{chap:na-practice} 章组织计算实验。最小二乘中的函数逼近依赖正交投影，插值则依赖节点值约束，两者采用不同的误差准则。

一学期课程可以把误差、求根、直接法、最小二乘、基本迭代、插值、数值积分和常微分方程作为主干，按专业需要选讲特征值、偏微分方程、傅里叶和蒙特卡洛专题。完整阅读适合配合持续的上机练习。正文中的例题展示方法与推导，章末习题包括证明、手算、反例和实验；附录给出解答或实验判据，实验中的预期现象不冒充实际运行结果。

\section{编写依据与资料使用范围}
课程范围参考三个可公开核验的来源。清华大学出版社对喻文健《数值分析与算法》第三版的介绍覆盖误差、求根、线性方程组、特征值、逼近、积分和常微分方程，并强调算法与应用\cite{yu2020}；北京大学公布的本科数值分析大纲包含函数逼近、数值微积分、非线性方程、常微分方程、快速傅里叶变换和蒙特卡洛方法\cite{pku-course}；MIT 18.330 的 2012 年公开课程突出离散化、收敛速度、傅里叶分析与数值实验，并说明其线性代数内容较少\cite{mit-syllabus,mit-notes}。

本书据此采用覆盖基础数值方法的主线，补充两类科学计算专题。上述资料用于课程定位、主题核对和延伸阅读，正文、例题与习题独立组织；本书并非上述学校的官方教材、原教材译本或逐章改写。由于只核验公开大纲与图书介绍，不把未查阅的纸质教材页码或具体证明列为引文依据。书中较深结果若不展开完整证明，会明确标示证明范围。

\section{记号与算法约定}
实数和复数分别记为 $\mathbb R$、$\mathbb C$，$\mathrm i^2=-1$。向量通常为列向量，$A^{\mathsf T}$ 表示转置，$A^*$ 表示共轭转置，$I$ 为适当阶数的单位矩阵。$\mathcal P_n$ 表示次数不超过 $n$ 的实多项式空间。$C^k[a,b]$ 表示在闭区间上具有连续至 $k$ 阶导数的函数；端点导数按单侧延拓理解。$\|f\|_\infty=\max_{x\in[a,b]}|f(x)|$。向量和矩阵范数在第~\ref{sec:na-norms} 节定义。

$x_*$ 或 $x^*$ 表示精确解，$x_k$ 表示第 $k$ 次迭代，$\widehat x$ 表示最终计算值，$e_k=x_k-x_*$ 表示有符号误差。网格宽度记为 $h$，时间步长记为 $\Delta t$ 或在明确的时间章节中记为 $h$，单位舍入误差记为 $u$。$E(h)=O(h^p)$ 表示存在与充分小的 $h>0$ 无关的常数 $C$，使 $|E(h)|\le Ch^p$；常数可依赖固定区间、函数或问题参数，若这些参数也变化，需另行说明。

算法描述在未注明时采用精确算术；实际程序需增加有限值检查、迭代上限和失败状态。给定容差不是给定真误差：只有附有误差界或满足相应条件的估计量，停止判据才具有误差含义。计算代价中的 $O(n^p)$ 以浮点加、减、乘、除为基本操作计数，忽略常数；函数评价昂贵时应另计其次数。

\chapter{误差、条件数与稳定性}
\label{chap:na-errors}

\section{误差的来源和衡量}
\begin{definition}[绝对误差与相对误差]
设 $x$ 是精确量，$\widehat x$ 是近似值。绝对误差（absolute error）为 $|\widehat x-x|$；当 $x\ne0$ 时，相对误差（relative error）为 $|\widehat x-x|/|x|$。向量情形用指定范数替代绝对值。若精确值为零，相对误差不定义，应使用绝对误差或事先约定的参考尺度。
\end{definition}

\keyterm{截断误差（truncation error）}来自用有限项级数、有限网格或有限差分替代无限或连续过程；\keyterm{舍入误差（rounding error）}来自有限精度数值表示和运算。迭代提前停止产生迭代误差。对一个固定模型，总误差可通过在精确解、精确离散解和实际计算解之间加入中间量后用三角不等式分解，但各项可能相互影响，不能假定它们总是独立随机量。

\begin{example}[相同绝对误差的不同含义]
长度 $1000$ 与 $0.001$ 都被测得绝对误差 $10^{-4}$ 时，相对误差分别为 $10^{-7}$ 与 $0.1$。前者足以支持许多相对精度要求，后者只有很粗略的相对精度。报告误差时应保留单位，或明确说明已经无量纲化。
\end{example}

\section{浮点模型与抵消}
\label{sec:na-floating}
规格化的基数为 $\beta$、有效位数为 $t$ 的浮点数（floating-point number）可写为 $\pm(d_0.d_1\cdots d_{t-1})_\beta\beta^e$，其中 $d_0\ne0$，指数 $e$ 位于有限范围。采用舍入到最近值且没有溢出、下溢影响时，单位舍入误差（unit roundoff）为 $u=\tfrac12\beta^{1-t}$。标准分析模型是
\begin{equation}
 \operatorname{fl}(a\circ b)=(a\circ b)(1+\delta),
 \qquad |\delta|\le u,\qquad \circ\in\{+,-,\times,\div\},
 \label{eq:na-rounding}
\end{equation}
除法要求 $b\ne0$，且精确结果处于该模型适用的范围。接近零的次正规数需要绝对误差模型；无穷大和非数值结果也不满足式~\eqref{eq:na-rounding}。常见二进制双精度格式的有效位数为 $53$，因此该约定下 $u=2^{-53}$；有的软件把 $1$ 与下一个较大浮点数的间距 $2^{-52}$ 称为机器精度，两种命名不能混用。

\begin{proposition}[连乘的误差合并]
若 $|\delta_j|\le u$ 且 $mu<1$，则
\[
 \prod_{j=1}^{m}(1+\delta_j)=1+\theta_m,
 \qquad |\theta_m|\le\gamma_m:=\frac{mu}{1-mu}.
\]
\end{proposition}
\begin{proof}
乘积位于 $(1-u)^m$ 与 $(1+u)^m$ 之间。Bernoulli 不等式给出 $(1-u)^m\ge1-mu$。展开上界并用 $\binom mj\le m^j$，得 $(1+u)^m\le\sum_{j=0}^\infty(mu)^j=(1-mu)^{-1}$。上下两侧的偏差均不超过 $\gamma_m$。
\end{proof}

把逐项顺序求和展开为每个输入乘上一串舍入因子，可得精确输入 $a_1,\ldots,a_n$ 的求和界
\[
 |\operatorname{fl}(\textstyle\sum_{j=1}^n a_j)-\sum_{j=1}^n a_j|
 \le\gamma_{n-1}\sum_{j=1}^n|a_j|,
 \qquad (n-1)u<1.
\]
因此正负大数相互抵消时，绝对误差可能仍小，而相对于很小的最终和的误差很大。相近数相减并非总在该次减法中产生巨大相对舍入；危险往往在于已有输入误差相对于相减结果被放大。

\begin{example}[等价公式的不同计算效果]
当 $x>0$ 很大时，$\sqrt{x^2+1}-x$ 的两个项很接近。恒等变形给出
\[
 \sqrt{x^2+1}-x=\frac1{\sqrt{x^2+1}+x}.
\]
右式避免了最后一步抵消。若 $x^2$ 可能溢出，还应使用带尺度处理的平方和开方例程，而不能只做有理化。另一例是小 $x$ 时的 $\exp(x)-1$，应使用针对该表达式设计的函数或适当级数，而非先计算接近 $1$ 的指数值再相减。
\end{example}

\section{范数与扰动的尺度}
\label{sec:na-norms}
\begin{definition}[向量范数与诱导矩阵范数]
向量范数（vector norm）满足正定性、绝对齐次性和三角不等式。对 $x\in\mathbb R^n$，常用
\[
 \|x\|_1=\sum_i|x_i|,\quad
 \|x\|_2=(\sum_i x_i^2)^{1/2},\quad
 \|x\|_\infty=\max_i|x_i|.
\]
对应的诱导矩阵范数（induced matrix norm）定义为 $\|A\|=\max_{x\ne0}\|Ax\|/\|x\|$。于是 $\|Ax\|\le\|A\|\|x\|$ 且 $\|AB\|\le\|A\|\|B\|$。
\end{definition}
具体地，$\|A\|_1=\max_j\sum_i|a_{ij}|$，$\|A\|_\infty=\max_i\sum_j|a_{ij}|$，$\|A\|_2=\sqrt{\lambda_{\max}(A^{\mathsf T}A)}$。后式来自对称矩阵正交谱分解与 Rayleigh 商。矩阵的 Frobenius 范数 $\|A\|_F=(\sum_{i,j}|a_{ij}|^2)^{1/2}$ 不是这里的诱导二范数，使用时应保留下标。

\section{问题的条件与算法的稳定性}
\label{sec:na-conditioning}
\begin{definition}[局部条件数]
若 $F$ 在 $x$ 附近可微，则数据扰动的绝对放大率由 $\|DF(x)\|$ 衡量。若 $x\ne0$ 且 $F(x)\ne0$，局部相对条件数（local relative condition number）定义为
\[
 \kappa_F(x)=\frac{\|DF(x)\|\|x\|}{\|F(x)\|}.
\]
标量情形为 $|xF'(x)/F(x)|$。它描述充分小、无结构限制的扰动的一阶最坏影响；允许扰动的集合改变时，条件数也可能改变。
\end{definition}

\begin{definition}[前向误差、后向误差与后向稳定]
求 $y=F(x)$ 得到 $\widehat y$ 时，$\|\widehat y-y\|$ 是前向误差（forward error）。若存在 $\Delta x$ 使 $\widehat y=F(x+\Delta x)$，则可以用最小的允许扰动大小衡量后向误差（backward error）。算法若使相对后向误差至多为 $Cu$，其中 $C$ 对给定维数受控且不随数据病态程度任意增大，称其在相应误差模型下后向稳定（backward stable）。
\end{definition}

小后向误差通过条件数转化为前向误差，因此常用关系是“相对前向误差约不超过条件数乘相对后向误差”，其中“约”指小扰动的一阶估计。病态问题即使使用稳定算法，仍可能有很大前向误差。稳定性也有分量意义、范数意义及保持结构的版本，不能只给算法贴上“稳定”标签而省略具体含义。

\begin{example}[相减问题的条件数]
令 $F(a,b)=a-b$，在输入采用无穷范数、输出采用绝对值时，$\|DF\|=2$，因而相对条件数为 $2\max(|a|,|b|)/|a-b|$。若改用逐分量相对扰动 $|\Delta a|\le\eta|a|$、$|\Delta b|\le\eta|b|$，放大率变为 $(|a|+|b|)/|a-b|$。两者都说明接近相等的数据相减容易放大已有误差，但数值表达式因扰动模型不同而不同。
\end{example}

\section{收敛阶与误差估计}
\begin{definition}[迭代的收敛阶]
设 $x_k\to x_*$，且误差从某步起非零。若存在 $p\ge1$ 与 $0<C<\infty$，满足
\[
 \lim_{k\to\infty}\frac{|x_{k+1}-x_*|}{|x_k-x_*|^p}=C,
\]
则称收敛阶为 $p$；$p=1$ 时还要求 $C<1$ 才称线性收敛，$p=2$ 时称二次收敛。若仅有 $|e_{k+1}|\le C|e_k|^p$，一般表述为至少具有相应局部误差上界，不据此断言极限恰等于一个正数。
\end{definition}

网格方法的 $p$ 阶精度通常指 $E(h)=O(h^p)$，不同于迭代次数意义下的收敛阶。若 $E(h)=Ch^p+o(h^p)$ 且 $C\ne0$，误差比可估计阶数；没有精确解时，用三个逐次减半的网格结果 $A_h,A_{h/2},A_{h/4}$ 构造
\[
 p_{\mathrm{obs}}=\log_2\frac{|A_h-A_{h/2}|}{|A_{h/2}-A_{h/4}|}.
\]
这一估计要求进入同一渐近区间且差值未被舍入主导，不能在差值为零时直接使用。

\section{习题}
\begin{enumerate}[label=\arabic*.]
\item\label{exr:na-error-rational} 将 $1-\cos x$ 改写成适于小 $x$ 的表达式，并解释改写解决的是哪种误差。
\item\label{exr:na-error-norm} 计算 $A=\begin{pmatrix}1&-2\\3&4\end{pmatrix}$ 的一范数、无穷范数和 Frobenius 范数。
\item\label{exr:na-error-condition} 求 $F(x)=\sqrt x$ 在 $x>0$ 时的绝对与相对条件数；说明为什么接近零时两个结论不矛盾。
\item\label{exr:na-error-balance} 设总误差上界为 $C_1h^2+C_2u/h$，其中 $C_1,C_2>0$。求使此上界最小的 $h$，并说明不能据此保证实际误差在该点取最小值。
\end{enumerate}

\chapter{非线性方程与不动点迭代}
\label{chap:na-roots}

\section{求根问题与二分法}
求根（root finding）要求给定 $f$ 后寻找 $x_*$ 使 $f(x_*)=0$。多个根、重根、定义域边界和不连续点都会改变算法行为。连续函数在区间两端异号保证至少有一根；二分法（bisection method）据此反复缩小包根区间。仅观察一张图或若干采样点不能替代连续异号条件。

\begin{theorem}[二分法的误差界]
设 $f\in C[a_0,b_0]$ 且 $f(a_0)f(b_0)<0$。每步取 $c_k=(a_k+b_k)/2$；若 $f(c_k)=0$ 则结束，否则保留两端异号的半区间作为 $[a_{k+1},b_{k+1}]$。在没有提前结束时，区间具有共同极限点 $x_*$，$f(x_*)=0$，且
\begin{equation}
 |c_k-x_*|\le\frac{b_0-a_0}{2^{k+1}}.
 \label{eq:na-bisection-bound}
\end{equation}
\end{theorem}
\begin{proof}
区间嵌套，长度为 $(b_0-a_0)2^{-k}$，故由实数完备性其交集是单点 $x_*$。若 $f(x_*)\ne0$，连续性使充分小邻域中的函数值同号，与各区间两端异号矛盾。中点与区间内任何点的距离不超过半个区间长度，得到式~\eqref{eq:na-bisection-bound}。
\end{proof}

输入应包括连续函数、初始包根区间、绝对位置容差 $\varepsilon_x>0$ 和迭代上限。停止于半区间长度不超过 $\varepsilon_x$，可得到位置保证；若只有函数值容差，还需要导数下界才能推出位置误差。初始化要单独处理端点恰为根。实现中比较符号优于直接计算 $f(a)f(b)$，因为乘积可能溢出或下溢；浮点中点若等于某个端点，应报告精度停滞。函数评价的符号必须可靠，含噪声评价不自动继承此定理。

\begin{example}[计算平方根的包根方法]
$f(x)=x^2-2$ 在 $[1,2]$ 连续且两端异号。取初始中点为第 $0$ 步，若要求误差不超过 $10^{-6}$，则需 $2^{-(k+1)}\le10^{-6}$，取 $k=19$ 即可，即至多评价 $20$ 个中点。该结论依赖区间收缩，与根附近导数大小无关。
\end{example}

\section{压缩映射与先验、后验误差}
满足 $g(x_*)=x_*$ 的点称为不动点（fixed point）；按 $x_{k+1}=g(x_k)$ 递推称为不动点迭代。压缩映射（contraction mapping）用一个小于 $1$ 的统一比例缩小任意两点的距离。

\begin{theorem}[区间上的压缩映射]
\label{thm:na-contraction}
设 $g:[a,b]\to[a,b]$，且存在 $0\le q<1$ 使 $|g(x)-g(y)|\le q|x-y|$。则 $g$ 有唯一不动点 $x_*$，任取 $x_0\in[a,b]$，迭代 $x_{k+1}=g(x_k)$ 收敛，并满足
\begin{align}
 |x_k-x_*|&\le q^k|x_0-x_*|\le\frac{q^k}{1-q}|x_1-x_0|,\label{eq:na-fixed-priori}\\
 |x_k-x_*|&\le\frac{q}{1-q}|x_k-x_{k-1}|\quad(k\ge1).\label{eq:na-fixed-posteriori}
\end{align}
\end{theorem}
\begin{proof}
连续两步差满足 $|x_{j+1}-x_j|\le q^j|x_1-x_0|$。几何级数尾和控制任意两项差，说明序列为 Cauchy 列；闭区间完备，极限属于区间。由 $g$ 的连续性，极限满足不动点等式。两个不动点距离不超过自身的 $q$ 倍，只能相同。把尾和从第 $k$ 步开始相加即得两个误差界；直接比较迭代点与不动点又得第一个不等式。
\end{proof}

若 $g$ 可微，$\max|g'|\le q<1$ 是压缩的充分条件，但还必须检查区间不变性 $g([a,b])\subseteq[a,b]$。局部条件 $|g'(x_*)|<1$ 在 $g'$ 连续时保证充分接近根的初值收敛，不保证任意初值收敛。

\begin{example}[求解 $x=\cos x$]
在 $[0,1]$ 上，$g(x)=\cos x$ 的值属于 $[\cos1,1]\subset[0,1]$，且 $|g'(x)|\le\sin1<1$。因此区间中有唯一不动点，任意区间内初值均收敛。式~\eqref{eq:na-fixed-posteriori} 将相邻迭代差转化为严格误差上界，而仅说“相邻两次差很小”没有这种保证。
\end{example}

\section{Newton 方法及其局部性质}
Newton 方法（Newton's method）以 $x_k$ 处的切线近似 $f$，令线性近似为零，得到
\begin{equation}
 x_{k+1}=x_k-\frac{f(x_k)}{f'(x_k)}.
 \label{eq:na-newton}
\end{equation}
这一步要求 $f'(x_k)\ne0$。每步需要函数与导数评价；比二分法更快的局部收敛必须与这些评价成本一起比较。

\begin{theorem}[Newton 法的局部二次误差界]
\label{thm:na-newton}
设 $f\in C^2$ 于根 $x_*$ 的某邻域，$f(x_*)=0$ 且 $f'(x_*)\ne0$。存在 $r>0$，使 $|x_0-x_*|\le r$ 时迭代有定义、留在该邻域并收敛，且
\[
 |x_{k+1}-x_*|\le \frac{M}{2m}|x_k-x_*|^2,
\]
其中可在该邻域取 $|f'|\ge m>0$、$|f''|\le M$。若误差不提前为零且 $f''(x_*)\ne0$，则二次收敛的渐近常数为 $|f''(x_*)/(2f'(x_*))|$。
\end{theorem}
\begin{proof}
在 $x_k$ 处对 $f(x_*)$ 作二阶 Taylor 展开，得
\[
 0=f(x_k)+f'(x_k)(x_*-x_k)+\tfrac12 f''(\xi_k)(x_*-x_k)^2.
\]
与式~\eqref{eq:na-newton} 相减可得 $e_{k+1}=f''(\xi_k)e_k^2/(2f'(x_k))$。由导数连续性取到上述上下界，再缩小 $r$ 使 $Mr/(2m)<1$，即保证误差缩小且迭代不离开邻域。取极限给出渐近常数。
\end{proof}

若 $f(x)=(x-x_*)^m q(x)$，整数 $m>1$ 且 $q(x_*)\ne0$，普通 Newton 迭代通常只有线性主项 $e_{k+1}=(1-1/m)e_k+O(e_k^2)$。已知重数时，改为 $x_{k+1}=x_k-mf(x_k)/f'(x_k)$ 可在适当光滑条件下消去线性项。根附近 $f'$ 很小也表示求根问题可能敏感，不能把所有困难归结为算法实现。

\begin{example}[局部结论不能当作全局结论]
对 $f(x)=x^3-2x+2$，从 $x_0=0$ 开始的 Newton 迭代满足 $x_1=1$、$x_2=0$，因而形成二周期。虽然该函数存在实根，算法仍不从这个初值收敛。一个实用做法是在已知包根区间内尝试 Newton 步，若候选点不满足预定的区间收缩要求，则退回二分步。若接受准则保证每次新区间长度至多为原来的固定比例 $\rho<1$，即可保留全局区间收敛保证。
\end{example}

\section{割线法、方程组与停止判据}
用差商替代导数得到\keyterm{割线法（secant method）}
\[
 x_{k+1}=x_k-f(x_k)\frac{x_k-x_{k-1}}{f(x_k)-f(x_{k-1})}.
\]
分母不能为零或因相消而过小。对充分光滑的简单根、足够接近的初值，在非退化情形下其收敛阶为 $(1+\sqrt5)/2$；直观上误差主项 $e_{k+1}\approx C e_ke_{k-1}$ 导出 $p^2=p+1$。这里是渐近关系的解释，并非局部收敛定理的完整证明。割线法每步新增一次函数评价，适合导数难求但函数评价可用的场合。

对 $F:\mathbb R^n\to\mathbb R^n$，Jacobian 矩阵为 $J_F(x)=(\partial F_i/\partial x_j)$。Newton 步通过
\[
 J_F(x_k)s_k=-F(x_k),\qquad x_{k+1}=x_k+s_k
\]
求得，不应为此显式形成矩阵逆。若 $J_F$ 在根附近 Lipschitz 连续且 $J_F(x_*)$ 可逆，则充分近的初值有局部二次误差界；理由是积分形式的 Taylor 余项为 $O(\|e_k\|^2)$，而附近逆矩阵范数有界。实际常采用阻尼 $x_{k+1}=x_k+\alpha_ks_k$，但任取 $0<\alpha_k<1$ 并不自动保证全局收敛。

停止时应同时观察步长和经过尺度化的残差，并报告达到容差、达到上限、导数失效等状态。标量情形若连接 $x$ 与唯一根 $x_*$ 的区间上 $|f'|\ge m>0$，中值定理给出 $|x-x_*|\le|f(x)|/m$。若 $f$ 整体乘以一个小常数，根不变而残差缩小，故未经尺度处理的固定残差阈值可能误导。

\section{习题}
\begin{enumerate}[label=\arabic*.]
\item\label{exr:na-root-bisect} 在 $[1,2]$ 用二分法求 $x^3-x-1=0$。证明根唯一，并求保证中点误差不超过 $10^{-8}$ 所需的中点评价次数上界。
\item\label{exr:na-root-fixed} 对 $x_{k+1}=(x_k+2/x_k)/2$，证明任意 $x_0>0$ 都收敛到 $\sqrt2$，并给出精确误差递推式。
\item\label{exr:na-root-multiple} 对 $f(x)=(x-1)^3$，比较普通 Newton 法与已知重数的修正方法。
\item\label{exr:na-root-system} 写出求解 $x^2+y^2=1$、$x-y=0$ 的 Newton 线性系统，并指出 Jacobian 何时奇异。
\end{enumerate}

\chapter{线性方程组的直接方法}
\label{chap:na-direct}

\section{线性系统、条件数与残差}
考虑 $Ax=b$，其中 $A\in\mathbb R^{n\times n}$ 非奇异。数学上 $x=A^{-1}b$，但计算上通常先分解 $A$ 再解三角系统。显式求逆要计算多余信息，且不比直接解方程自动更准确。

\begin{definition}[矩阵条件数]
在同一种诱导范数下，非奇异矩阵的条件数为 $\kappa(A)=\|A\|\|A^{-1}\|$。由于 $\|I\|=1$，有 $\kappa(A)\ge1$。二范数条件数等于最大奇异值与最小奇异值之比。
\end{definition}
若只扰动右端项，则 $A\Delta x=\Delta b$。当 $b\ne0$ 时，结合 $\|b\|\le\|A\|\|x\|$ 得
\begin{equation}
 \frac{\|\Delta x\|}{\|x\|}\le\kappa(A)\frac{\|\Delta b\|}{\|b\|}.
 \label{eq:na-rhs-perturbation}
\end{equation}

\begin{theorem}[系数与右端项同时扰动]
\label{thm:na-linear-perturbation}
设 $(A+\Delta A)(x+\Delta x)=b+\Delta b$，$Ax=b\ne0$，并令 $\eta_A=\|\Delta A\|/\|A\|$、$\eta_b=\|\Delta b\|/\|b\|$。若 $\kappa(A)\eta_A<1$，则 $A+\Delta A$ 非奇异，而且
\[
 \frac{\|\Delta x\|}{\|x\|}\le
 \frac{\kappa(A)}{1-\kappa(A)\eta_A}(\eta_A+\eta_b).
\]
\end{theorem}
\begin{proof}
令 $E=A^{-1}\Delta A$。由 $\|E\|<1$，矩阵级数 $\sum_{j=0}^{\infty}(-E)^j$ 在范数下绝对收敛，乘以 $I+E$ 后为 $I$，故 $\|(I+E)^{-1}\|\le(1-\|E\|)^{-1}$。两方程相减给出 $(I+E)\Delta x=A^{-1}\Delta b-Ex$。取范数，再用 $\|b\|/\|x\|\le\|A\|$ 即得结论。
\end{proof}

对近似解 $\widehat x$ 定义残差 $r=b-A\widehat x$。因为 $A(x-\widehat x)=r$，有
\[
 \frac{\|x-\widehat x\|}{\|x\|}\le\kappa(A)\frac{\|r\|}{\|b\|}.
\]
若只允许右端项改变，$\widehat x$ 是 $A\widehat x=b-r$ 的精确解，因此 $\|r\|/\|b\|$ 是相应相对后向误差。常见尺度化指标 $\|r\|/(\|A\|\|\widehat x\|+\|b\|)$ 还考虑了系数与解的尺度；报告时应注明范数，不能把它直接称为前向误差。

\begin{example}[病态系统中的小残差]
令 $A=\operatorname{diag}(1,\varepsilon)$、$b=(1,\varepsilon)^{\mathsf T}$，其中 $0<\varepsilon<1$。精确解为 $(1,1)^{\mathsf T}$，近似解 $(1,0)^{\mathsf T}$ 的无穷范数相对误差为 $1$，相对残差为 $\varepsilon$，而 $\kappa_\infty(A)=1/\varepsilon$。残差界恰能解释误差为何没有随残差一起变小。
\end{example}

\section{消元、三角求解与选主元}
\keyterm{高斯消元（Gaussian elimination）}逐列消去主对角线下方元素。第 $k$ 步用非零主元 $u_{kk}$ 形成乘子 $l_{ik}=a_{ik}^{(k)}/u_{kk}$，再更新
\[
 a_{ij}^{(k+1)}=a_{ij}^{(k)}-l_{ik}a_{kj}^{(k)},\qquad i,j>k.
\]
若各步主元非零，消元等价于 $A=LU$，其中 $L$ 是单位下三角矩阵，$U$ 是上三角矩阵。先解 $Ly=b$，再解 $Ux=y$。回代公式为
\[
 x_i=\frac{y_i-\sum_{j=i+1}^{n}u_{ij}x_j}{u_{ii}},\qquad i=n,n-1,\ldots,1.
\]
前代与回代各需 $O(n^2)$ 次运算；稠密消元约需 $2n^3/3$ 次浮点运算。对同一矩阵的多个右端项应复用分解。

\begin{proposition}[无换行 LU 分解的充分条件]
若 $A$ 的所有顺序主子式均非零，则存在唯一的单位下三角 $L$ 与非奇异上三角 $U$，使 $A=LU$。
\end{proposition}
\begin{proof}
按阶数归纳。第一主元非零，可消去第一列，其余块为 Schur 补。该块前 $k$ 阶主子式等于原矩阵前 $k+1$ 阶主子式除以第一主元，因而非零，可继续归纳。唯一性来自 $L_2^{-1}L_1=U_2U_1^{-1}$：该矩阵同时为单位下三角与上三角，只能等于 $I$。
\end{proof}

\keyterm{部分选主元（partial pivoting）}在当前列的剩余元素中选绝对值最大的作为主元，必要时交换行，得到
\begin{equation}
 PA=LU,
 \label{eq:na-plu}
\end{equation}
其中 $P$ 是置换矩阵。交换行时，已经存入 $L$ 的前几列乘子也必须相应交换。输入为方阵 $A$ 与右端项，输出为解或奇异、近奇异等状态；选主元为零说明精确算术下矩阵奇异，浮点下则还应结合矩阵尺度判断过小主元。

\begin{example}[小主元引起的中间量放大]
对 $A=\begin{pmatrix}\varepsilon&1\\1&1\end{pmatrix}$，直接以 $\varepsilon$ 为主元产生乘子 $1/\varepsilon$。当 $\varepsilon$ 很小时，中间量很大；先交换两行则乘子为 $\varepsilon$。原矩阵在 $\varepsilon\to0$ 时仍趋于非奇异矩阵，这说明坏消元步骤可以出现在良态问题上。
\end{example}

部分选主元保证乘子绝对值不超过 $1$，但并不禁止 $U$ 的元素增长。标准误差分析在通常浮点模型下给出与 $u|L||U|$ 同阶的分量后向误差界，常数依赖维数。由此可知其稳定性还受消元增长因子影响；不能断言它对一切矩阵都有与数据无关的小误差常数。工程中它通常可靠，严格分析仍需保留增长因素。

\section{对称正定系统与 Cholesky 分解}
\begin{definition}[对称正定矩阵]
实对称矩阵 $A$ 若对所有 $x\ne0$ 满足 $x^{\mathsf T}Ax>0$，称为对称正定矩阵（symmetric positive definite matrix，SPD）。
\end{definition}
\begin{theorem}[Cholesky 分解]
若 $A$ 对称正定，则存在唯一的下三角矩阵 $L$，其对角元为正，使 $A=LL^{\mathsf T}$。
\end{theorem}
\begin{proof}
写 $A=\begin{pmatrix}a&v^{\mathsf T}\\v&B\end{pmatrix}$，则 $a>0$。对 $y\ne0$，令 $z=(-v^{\mathsf T}y/a,y^{\mathsf T})^{\mathsf T}$，有 $y^{\mathsf T}(B-vv^{\mathsf T}/a)y=z^{\mathsf T}Az>0$，故 Schur 补仍正定。取 $l_{11}=\sqrt a$、首列下部为 $v/\sqrt a$，对子块递归即可存在。逐列公式与正对角元约定又确定了唯一性。
\end{proof}
具体公式为
\[
 l_{jj}=\left(a_{jj}-\sum_{k<j}l_{jk}^2\right)^{1/2},\qquad
 l_{ij}=\frac{a_{ij}-\sum_{k<j}l_{ik}l_{jk}}{l_{jj}},\quad i>j.
\]
稠密计算约需 $n^3/3$ 次浮点运算。若实际计算出现非正开方量，应报告分解失败并检查正定性、尺度和舍入影响；不应把负值随意取绝对值来强行完成分解。

\section{结构化系统与迭代改进}
三对角系统只在主对角线及两侧相邻对角线上有非零元。无换行消元若每个主元非零，只需 $O(n)$ 存储和运算；一般三对角矩阵仍可能出现零主元，正定三对角情形则可使用其结构化 Cholesky 分解。稀疏矩阵的消元还会产生原先不存在的非零元，称为\keyterm{填充（fill-in）}，所以“输入稀疏”不等于“分解仍同样稀疏”。

\keyterm{迭代改进（iterative refinement）}先求近似解 $\widehat x$，再计算 $r=b-A\widehat x$，用已有分解近似求解 $Ad=r$，更新 $\widehat x\leftarrow\widehat x+d$。精确算术下若修正方程解得精确，一次修正即可恢复真解；浮点效果取决于分解误差、条件数与残差评价精度。以较高精度算残差常有帮助，但并非对任意病态问题都保证恢复全部有效数字。

\section{习题}
\begin{enumerate}[label=\arabic*.]
\item\label{exr:na-direct-lu} 求 $A=\begin{pmatrix}2&1\\4&3\end{pmatrix}$ 的无换行 LU 分解，并解 $Ax=(1,3)^{\mathsf T}$。
\item\label{exr:na-direct-cholesky} 求 $A=\begin{pmatrix}4&2\\2&3\end{pmatrix}$ 的 Cholesky 分解。
\item\label{exr:na-direct-condition} 证明 $\kappa_2(Q)=1$ 对任意实正交矩阵成立，并求 $\operatorname{diag}(1,10^{-6})$ 的二范数条件数。
\item\label{exr:na-direct-residual} 已知 $\kappa(A)=10^5$，某近似解的相对残差为 $10^{-10}$。能保证怎样的相对前向误差上界？能否由此断言误差恰为该上界？
\end{enumerate}

\chapter{最小二乘、正交分解与函数逼近}
\label{chap:na-ls}

\section{超定方程与正交投影}
测量数据往往给出比未知数更多的方程，噪声使它们无法同时满足。给定 $A\in\mathbb R^{m\times n}$、$b\in\mathbb R^m$ 且 $m\ge n$，\keyterm{线性最小二乘（linear least squares）}求
\begin{equation}
 \min_{x\in\mathbb R^n}\|Ax-b\|_2^2.
 \label{eq:na-ls-objective}
\end{equation}
这里优化变量为 $x$，数据为 $A,b$，目标是残差平方和，而非插值所有数据。

\begin{theorem}[正规方程与唯一性]
\label{thm:na-normal-equations}
向量 $x_*$ 是式~\eqref{eq:na-ls-objective} 的极小点，当且仅当
\begin{equation}
 A^{\mathsf T}Ax_*=A^{\mathsf T}b.
 \label{eq:na-normal-equations}
\end{equation}
式~\eqref{eq:na-normal-equations} 称为正规方程（normal equations）。最小二乘解总存在；若 $A$ 列满秩则唯一，此时 $A^{\mathsf T}A$ 对称正定。
\end{theorem}
\begin{proof}
最优点处沿任意方向 $z$ 的一阶导数为 $2z^{\mathsf T}A^{\mathsf T}(Ax_*-b)=0$，故正规方程必要。反之，若 $r_*=b-Ax_*$ 满足 $A^{\mathsf T}r_*=0$，则
\[
 \|A(x_*+z)-b\|_2^2=\|r_*\|_2^2+\|Az\|_2^2,
\]
故它为极小点。有限维列空间是闭子空间，对 $b$ 的正交投影存在，投影向量可表示为 $Ax_*$，从而解存在。列满秩时非零 $z$ 满足 $Az\ne0$，故极小点唯一且正规矩阵正定。
\end{proof}

正规方程是理论刻画，不代表数值上总应显式形成它。若 $A$ 列满秩，则 $\kappa_2(A^{\mathsf T}A)=\kappa_2(A)^2$，因此形成正规方程会放大条件方面的困难。这里矩形矩阵的二范数条件数按最大、最小奇异值之比定义。

\begin{example}[直线拟合]
对数据 $(0,1),(1,2),(2,2)$，拟合 $p(t)=\alpha+\beta t$。设计矩阵与正规方程为
\[
 A=\begin{pmatrix}1&0\\1&1\\1&2\end{pmatrix},\quad
 A^{\mathsf T}A=\begin{pmatrix}3&3\\3&5\end{pmatrix},\quad
 A^{\mathsf T}b=\begin{pmatrix}5\\6\end{pmatrix}.
\]
解得 $\alpha=7/6$、$\beta=1/2$。按 $b-Ax$ 约定，残差为 $(-1/6,1/3,-1/6)^{\mathsf T}$，与设计矩阵两列都正交，残差平方和为 $1/6$。拟合曲线没有穿过每一个点，这是目标函数允许的行为。
\end{example}

\section{QR 分解与 Householder 变换}
列满秩矩阵可写为约化 QR 分解 $A=Q_1R$，其中 $Q_1\in\mathbb R^{m\times n}$、$Q_1^{\mathsf T}Q_1=I$，$R\in\mathbb R^{n\times n}$ 为非奇异上三角矩阵。将 $Q_1$ 扩充为正交矩阵 $Q=(Q_1,Q_2)$，利用正交变换保二范数，有
\[
 \|Ax-b\|_2^2=\|Rx-Q_1^{\mathsf T}b\|_2^2+\|Q_2^{\mathsf T}b\|_2^2.
\]
所以只需求解 $Rx=Q_1^{\mathsf T}b$。当 $m=n$ 时第二项为空项，视为零。

\keyterm{Householder 反射（Householder reflection）}定义为 $H=I-2vv^{\mathsf T}/(v^{\mathsf T}v)$，其中 $v\ne0$。直接相乘可验证 $H^{\mathsf T}=H$、$H^2=I$。给定非零向量 $z$，取 $\alpha=-\operatorname{sign}(z_1)\|z\|_2$，约定 $\operatorname{sign}(0)=1$，再取 $v=z-\alpha e_1$，即可使 $Hz=\alpha e_1$。符号选择避免 $v$ 首分量由两个接近的大数相减得到。

逐列把反射作用于剩余子矩阵即可得到 QR 分解，无须显式存储每个完整 $H$；更新可写为矩阵与向量的运算。稠密 $m\times n$ 矩阵且 $m\ge n$ 的代价约为 $2mn^2-2n^3/3$。Householder QR 在标准浮点条件下有范数意义的后向稳定性。经典 Gram--Schmidt 正交化在接近线性相关时可能丢失正交性，修正版本和重新正交化能改善这一现象；这些算法得到相同精确分解的思想，却有不同有限精度表现。

\section{奇异值分解与正则化}
\begin{theorem}[奇异值分解的基本形式]
任意 $A\in\mathbb R^{m\times n}$ 可写为 $A=U\Sigma V^{\mathsf T}$，其中 $U,V$ 为正交矩阵，$\Sigma$ 为矩形对角矩阵，非负对角元 $\sigma_1\ge\cdots\ge\sigma_r>0$，其余为零，$r=\operatorname{rank}(A)$。这些正数称为奇异值（singular values）。
\end{theorem}
\begin{proof}[证明要点]
对称半正定矩阵 $A^{\mathsf T}A$ 有正交特征向量组 $v_i$，非负特征值可写为 $\sigma_i^2$。对 $\sigma_i>0$ 定义 $u_i=Av_i/\sigma_i$，可验证这些向量正交归一。把它们补成 $\mathbb R^m$ 的正交基，并对零特征值对应的核空间选取正交基，即得分解。这里使用线性代数中的实对称谱定理。
\end{proof}

由奇异向量展开，最小二乘解中二范数最小者为
\[
 x^\dagger=\sum_{i=1}^r\frac{u_i^{\mathsf T}b}{\sigma_i}v_i.
\]
在 $\ker A$ 方向增加任意向量不会改变残差，却增大该最小范数解的范数。很小的 $\sigma_i$ 会放大数据中对应方向的扰动，形成逆问题中的不稳定性。

\keyterm{Tikhonov 正则化（Tikhonov regularization）}在本书采用参数平方约定：给定 $\lambda>0$，求
\[
 \min_x\bigl(\|Ax-b\|_2^2+\lambda^2\|x\|_2^2\bigr),\qquad
 x_\lambda=\sum_{i=1}^r\frac{\sigma_i}{\sigma_i^2+\lambda^2}(u_i^{\mathsf T}b)v_i.
\]
新的目标有唯一解，因为 $A^{\mathsf T}A+\lambda^2I$ 正定。正则化通过减弱小奇异值方向引入偏差并降低扰动放大；选择 $\lambda$ 需要噪声水平、验证数据或其他先验信息，不能仅以训练残差最小来决定。

\section{连续最小二乘与正交多项式}
设 $w$ 在 $[a,b]$ 上连续且严格为正，定义 $\langle f,g\rangle_w=\int_a^b f(x)g(x)w(x)\,\mathrm dx$。给定连续函数 $f$ 与有限维子空间 $V=\operatorname{span}\{\phi_0,\ldots,\phi_n\}$，最佳平方逼近求 $p\in V$ 使 $\|f-p\|_w^2=\langle f-p,f-p\rangle_w$ 最小。若基函数线性无关，则 Gram 矩阵 $G_{ij}=\langle\phi_j,\phi_i\rangle_w$ 正定，存在唯一解，且 $f-p$ 与 $V$ 正交。证明与定理~\ref{thm:na-normal-equations} 相同，只需把欧氏内积替换为积分内积。

当基函数两两正交时，系数直接为 $c_j=\langle f,\phi_j\rangle_w/\langle\phi_j,\phi_j\rangle_w$。在 $[-1,1]$、$w=1$ 下，前几个 Legendre 多项式是 $1,x,(3x^2-1)/2$。选择正交基改善系数求解的表示条件，但连续最佳平方逼近仍不同于一致范数下的最佳逼近，也不同于有限样本拟合。

\begin{example}[平方函数的最佳一次平方逼近]
在 $[-1,1]$ 上以权 $1$ 逼近 $f(x)=x^2$。由于奇偶性，$x$ 的系数为零；常数项为 $\int_{-1}^1x^2\,\mathrm dx/\int_{-1}^1 1\,\mathrm dx=1/3$。故 $p(x)=1/3$，平方误差为 $\int_{-1}^1(x^2-1/3)^2\,\mathrm dx=8/45$。
\end{example}

\section{习题}
\begin{enumerate}[label=\arabic*.]
\item\label{exr:na-ls-constant} 用常数拟合 $b_1,\ldots,b_m$，证明最小二乘解为样本平均数；给出正权 $w_i$ 下的加权结果。
\item\label{exr:na-ls-householder} 证明 $H=I-2vv^{\mathsf T}/(v^{\mathsf T}v)$ 保持二范数，且将 $v$ 映为 $-v$。
\item\label{exr:na-ls-rank} 求 $A=\begin{pmatrix}1&1\\2&2\end{pmatrix}$、$b=(1,2)^{\mathsf T}$ 的全部最小二乘解及最小范数解。
\item\label{exr:na-ls-projection} 求 $x^3$ 在 $[-1,1]$、权 $1$ 下的最佳一次平方逼近及平方误差。
\end{enumerate}

\chapter{线性方程组的迭代方法}
\label{chap:na-iterative}

\section{矩阵分裂与收敛的谱条件}
对于大型稀疏系统，直接分解的填充和存储代价可能很大。迭代法利用矩阵乘向量和简单子问题逐步逼近解。设 $A=M-N$，$M$ 可逆，令
\begin{equation}
 x_{k+1}=M^{-1}Nx_k+M^{-1}b=:Bx_k+c.
 \label{eq:na-stationary}
\end{equation}
实现中求解关于 $M$ 的方程，不必显式形成 $M^{-1}$。误差满足 $e_{k+1}=Be_k$。

\begin{theorem}[定常迭代的充要条件]
对任意初值，式~\eqref{eq:na-stationary} 收敛到唯一不动点，当且仅当 $\rho(B)<1$，其中谱半径 $\rho(B)=\max_{\lambda\in\sigma(B)}|\lambda|$，$\sigma(B)$ 为复特征值集合。
\end{theorem}
\begin{proof}
若对任意初始误差有 $B^ke_0\to0$，则对每个复特征向量亦有 $\lambda^kv\to0$，故 $|\lambda|<1$。反之，对复 Jordan 标准形，每个块的 $k$ 次幂由有限个 $k$ 的多项式乘 $\lambda^{k-j}$ 构成；当 $|\lambda|<1$ 时均趋于零，零特征值块最终为零。因此 $B^k\to0$，且 $I-B$ 可逆，不动点唯一。
\end{proof}

容易检查的充分条件为某诱导范数下 $\|B\|<1$。它不是必要条件；$\rho(B)<1$ 时某些范数下仍可能 $\|B\|>1$，初期误差也不必单调下降。

\section{Jacobi、Gauss--Seidel 与松弛}
写 $A=D+L+U$，其中 $D$ 为对角部分，$L,U$ 分别是严格下、上三角部分。本章采用此加号约定。若所有 $a_{ii}\ne0$，Jacobi 迭代为
\[
 x_i^{(k+1)}=\frac{b_i-\sum_{j\ne i}a_{ij}x_j^{(k)}}{a_{ii}},
 \qquad B_J=-D^{-1}(L+U).
\]
Gauss--Seidel 迭代在一轮内立即使用已更新分量：
\[
 x_i^{(k+1)}=\frac{b_i-\sum_{j<i}a_{ij}x_j^{(k+1)}-\sum_{j>i}a_{ij}x_j^{(k)}}{a_{ii}},
 \qquad B_G=-(D+L)^{-1}U.
\]
若 $A$ 按行严格对角占优，即 $|a_{ii}|>\sum_{j\ne i}|a_{ij}|$，则 Jacobi 迭代在无穷范数下为压缩映射。Gauss--Seidel 也收敛：令旧误差最大值为 $E$，$q=\max_i\sum_{j\ne i}|a_{ij}|/|a_{ii}|<1$；按行归纳，新误差每个分量都不超过 $qE$，因为本轮已更新的误差至多为 $E$，剩余旧误差也至多为 $E$。

逐分量把新 Gauss--Seidel 值与旧值以权 $\omega$ 混合，得到逐次超松弛法（successive over-relaxation，SOR）：
\[
 (D+\omega L)x_{k+1}=\omega b-\bigl(\omega U+(\omega-1)D\bigr)x_k.
\]
对称正定矩阵在 $0<\omega<2$ 时 SOR 收敛；这一标准结果可由逐坐标二次能量下降证明，本书不展开其完整谱分析。对一般矩阵，$0<\omega<2$ 本身不是充分条件，也不存在与矩阵无关的最佳参数。

\begin{example}[一个可直接比较的系统]
对 $A=\begin{pmatrix}4&1\\1&3\end{pmatrix}$，Jacobi 矩阵为 $\begin{pmatrix}0&-1/4\\-1/3&0\end{pmatrix}$，谱半径为 $1/\sqrt{12}$；Gauss--Seidel 矩阵为 $\begin{pmatrix}0&-1/4\\0&1/12\end{pmatrix}$，谱半径为 $1/12$。二者都收敛，且后者在此例具有更小的渐近因子；该比较不能无条件推广到所有矩阵。
\end{example}

\section{最速下降与共轭梯度}
对 SPD 矩阵 $A$，求解 $Ax=b$ 等价于最小化
\[
 \phi(x)=\tfrac12x^{\mathsf T}Ax-b^{\mathsf T}x,
 \qquad \phi(x)-\phi(x_*)=\tfrac12\|x-x_*\|_A^2,
\]
其中 $\|z\|_A=(z^{\mathsf T}Az)^{1/2}$ 为能量范数。残差 $r_k=b-Ax_k$ 是负梯度。沿 $r_k$ 精确线搜索得最速下降步
\[
 \alpha_k=\frac{r_k^{\mathsf T}r_k}{r_k^{\mathsf T}Ar_k},\qquad
 x_{k+1}=x_k+\alpha_kr_k.
\]
当残差非零时分母为正。条件数很大时不同方向曲率差异显著，最速下降可能反复折返，收敛缓慢。

\keyterm{共轭梯度法（conjugate gradient method，CG）}使用互相 $A$ 正交的搜索方向。给定 $x_0$，初始化 $r_0=b-Ax_0$、$p_0=r_0$；若 $r_0=0$ 直接返回。其精确算术递推为
\begin{align}
 \alpha_k&=\frac{r_k^{\mathsf T}r_k}{p_k^{\mathsf T}Ap_k},&
 x_{k+1}&=x_k+\alpha_kp_k,\notag\\
 r_{k+1}&=r_k-\alpha_kAp_k,&
 \beta_k&=\frac{r_{k+1}^{\mathsf T}r_{k+1}}{r_k^{\mathsf T}r_k},\notag\\
 p_{k+1}&=r_{k+1}+\beta_kp_k.
 \label{eq:na-cg}
\end{align}
若新残差为零，应在计算下一步前停止。每轮一次矩阵向量乘法及 $O(n)$ 次向量操作；对稀疏矩阵，乘法代价与非零元数同阶。

\begin{theorem}[CG 的投影性质与误差界]
\label{thm:na-cg}
在精确算术下，SPD 系统的 CG 在尚未到达解时产生 $A$ 共轭的非零搜索方向。第 $k$ 步在 $x_0+\mathcal K_k$ 中最小化能量误差，其中
\[
 \mathcal K_k=\operatorname{span}\{r_0,Ar_0,\ldots,A^{k-1}r_0\}.
\]
故它至多 $n$ 步终止。若 $\kappa=\kappa_2(A)>1$，则
\[
 \|x_k-x_*\|_A\le
 2\left(\frac{\sqrt\kappa-1}{\sqrt\kappa+1}\right)^k\|x_0-x_*\|_A.
\]
当 $\kappa=1$ 时非精确初值一步求解。
\end{theorem}
\begin{proof}[证明思路与误差界的依据]
将式~\eqref{eq:na-cg} 代入内积关系，可归纳得到不同残差正交、不同方向 $A$ 正交，以及两组向量生成相应 Krylov 空间。沿共轭方向的精确极小化不会破坏先前方向上的最优性，因此 $r_k$ 与 $\mathcal K_k$ 正交，这正是仿射空间中的极小化条件。非零共轭方向线性无关，故至多有 $n$ 个。由极小化性质，任取次数不超过 $k$ 且 $p(0)=1$ 的多项式，均有
\[
 \|x_k-x_*\|_A\le\max_{\lambda\in\sigma(A)}|p(\lambda)|\,\|x_0-x_*\|_A.
\]
把第~\ref{sec:na-chebyshev} 节的 Chebyshev 多项式从 $[-1,1]$ 映到 $[\lambda_{\min},\lambda_{\max}]$，并在零点归一化，即得所述上界。这里交代投影与多项式界的关键逻辑，递推正交关系的逐项归纳略去。
\end{proof}

有限精度中，共轭性与正交性可能逐渐丢失，不能把“至多 $n$ 步”当作实际软件的保证。递推残差也可能偏离真实残差，应在必要时重新计算 $b-Ax_k$。CG 的正定假设不可随意取消；非对称系统通常需要其他 Krylov 方法。

\section{预条件与可解释的停止}
\keyterm{预条件（preconditioning）}通过容易求解的近似矩阵改善迭代问题。对 SPD 预条件矩阵 $M=CC^{\mathsf T}$，等价变换为
\[
 C^{-1}AC^{-\mathsf T}y=C^{-1}b,\qquad x=C^{-\mathsf T}y.
\]
变换后的矩阵仍 SPD，适用于 CG。实现中用 $Mz=r$ 的求解代替显式形成 $C^{-1}$。Jacobi 预条件取 $M=D$，对 SPD 系统对角元为正。有效性取决于谱分布及每步成本，不仅是某个条件数数字。

迭代停止常用 $\|r_k\|\le\varepsilon_{\mathrm{abs}}+\varepsilon_{\mathrm{rel}}\|b\|$，配合最大迭代次数和停滞检测。该条件给出残差保证，前向误差需结合第~\ref{chap:na-direct} 章的条件数界。零右端项不能只用相对阈值；预条件残差和原问题残差也应明确区分。

\section{习题}
\begin{enumerate}[label=\arabic*.]
\item\label{exr:na-iter-jacobi} 对 $A=\begin{pmatrix}2&1\\1&2\end{pmatrix}$ 写出 Jacobi 和 Gauss--Seidel 迭代矩阵，求谱半径。
\item\label{exr:na-iter-nonnormal} 令 $B=\begin{pmatrix}1/2&10\\0&1/2\end{pmatrix}$。计算 $B^k$，说明 $\rho(B)<1$ 不保证每步误差范数下降。
\item\label{exr:na-iter-cg} 对 $A=\operatorname{diag}(1,2)$、$b=(1,1)^{\mathsf T}$、$x_0=0$，手算 CG 的前两步。
\item\label{exr:na-iter-energy} 证明 SPD 系统中 $\|x-x_*\|_A^2=r^{\mathsf T}A^{-1}r$，其中 $r=b-Ax$。
\end{enumerate}

\chapter{矩阵特征值的数值计算}
\label{chap:na-eigen}

\section{特征对、归一化与残差}
求特征对要求 $Av=\lambda v$ 且 $v\ne0$。特征向量可以任意非零倍缩放，因此通常约定 $\|v\|_2=1$，比较时还要考虑符号或复相位。近似特征对的残差为 $r=Av-\mu v$。单个特征向量在重复或聚集特征值附近可能很敏感，此时不变子空间比特定基向量更有意义。

\begin{proposition}[对称矩阵的残差界]
设 $A$ 实对称，$\|v\|_2=1$，$\mu\in\mathbb R$。则至少有一个特征值满足 $|\lambda-\mu|\le\|Av-\mu v\|_2$。
\end{proposition}
\begin{proof}
取正交特征基 $q_i$，写 $v=\sum_i c_iq_i$，$\sum_i c_i^2=1$。则残差平方为 $\sum_i(\lambda_i-\mu)^2c_i^2$，不小于 $\min_i|\lambda_i-\mu|^2$，开平方即得结论。
\end{proof}
这一定理给出某个特征值接近的保证，并不直接保证 $v$ 接近某个唯一特征向量；后者还需要谱间隙。

\section{幂法与反幂法}
幂法输入非零 $v_0$，每步计算 $w_k=Av_k$，若 $w_k\ne0$ 则令 $v_{k+1}=w_k/\|w_k\|_2$。若 $A$ 可对角化，存在唯一模最大的特征值 $|\lambda_1|>|\lambda_2|\ge\cdots$，且 $v_0$ 在其特征向量方向的展开系数不为零，则归一化方向趋近于该主特征方向，误差的主导速率由 $|\lambda_2/\lambda_1|^k$ 控制。若主特征值为负，向量符号可以交替；严谨的收敛对象是方向。证明直接来自 $A^kv_0=\sum_i c_i\lambda_i^kv_i$ 并除以 $\lambda_1^k$。

对实对称矩阵，可用\keyterm{Rayleigh 商（Rayleigh quotient）}
\[
 \mu(v)=\frac{v^{\mathsf T}Av}{v^{\mathsf T}v}
\]
估计特征值。若单位向量与某特征向量夹角为 $\theta$，正交展开显示该商的误差为 $O(\sin^2\theta)$，常数受其他特征值距离控制，故其精度常优于方向误差的一阶量。

\begin{example}[缺失主方向的初值]
对 $A=\operatorname{diag}(3,1)$，初值 $(1,1)^{\mathsf T}$ 产生分量比 $3^k:1$，归一化方向趋于第一坐标轴；初值 $(0,1)^{\mathsf T}$ 却永远留在第二坐标轴。主特征值存在且简单并不足以保证任意非零初值都成功。
\end{example}

为寻找最接近给定移位 $\sigma$ 的特征值，反幂法每步解 $(A-\sigma I)w_k=v_k$ 再归一化。若 $\sigma$ 不是特征值，则逆矩阵的特征值为 $(\lambda_i-\sigma)^{-1}$；唯一最近且被初值激活的方向成为主方向。固定移位可复用矩阵分解；过于接近特征值时求解系统病态，应同时监测原特征问题的残差。

\section{QR 迭代与实际计算范围}
不带移位的 QR 迭代从 $A_0=A$ 出发，每步作 $A_k=Q_kR_k$，再令 $A_{k+1}=R_kQ_k=Q_k^{\mathsf T}A_kQ_k$。每一步是正交相似变换，保持特征值。带移位时分解 $A_k-\sigma_kI=Q_kR_k$，再取 $A_{k+1}=R_kQ_k+\sigma_kI$，同样保持相似性。

实际稠密算法先把一般矩阵化为 Hessenberg 形，对称矩阵化为三对角形，再使用移位与消去已收敛块的策略。一般实矩阵可能含复共轭特征值，因此实 Schur 形允许 $2\times2$ 对角块，不能要求所有实矩阵都迭代到实对角阵。本节说明算法结构，不给出一般带移位 QR 算法的完整收敛理论；课程实验宜先选实对称矩阵，并以残差和正交性核验结果。

\section{习题}
\begin{enumerate}[label=\arabic*.]
\item\label{exr:na-eigen-power} 对 $\operatorname{diag}(4,2,1)$ 及初值 $(1,1,1)^{\mathsf T}$ 写出第 $k$ 次归一化幂迭代向量，并指出主导方向误差比例。
\item\label{exr:na-eigen-rayleigh} 对 $A=\operatorname{diag}(\lambda_1,\lambda_2)$、$v=(\cos\theta,\sin\theta)^{\mathsf T}$，求 Rayleigh 商误差的精确表达式。
\item\label{exr:na-eigen-shift} 对 $A=\operatorname{diag}(1,3,8)$，移位 $\sigma=2.8$ 的反幂法通常趋向哪一个特征方向？注明初值条件。
\item\label{exr:na-eigen-rotation} 解释实旋转矩阵 $\begin{pmatrix}0&-1\\1&0\end{pmatrix}$ 为什么不能正交相似于实对角阵。
\end{enumerate}

\chapter{多项式插值与分段逼近}
\label{chap:na-interpolation}

\section{插值问题与 Lagrange 形式}
给定互不相同的节点 $x_0,\ldots,x_n$ 和数据 $y_0,\ldots,y_n$，插值要求寻找 $p_n\in\mathcal P_n$ 满足 $p_n(x_j)=y_j$。它把离散信息延拓为函数表示，可用于评价、求导和积分；这些后续运算对误差的放大方式并不相同。

\begin{theorem}[多项式插值的存在唯一性]
上述插值问题有唯一解，可写为
\begin{equation}
 p_n(x)=\sum_{j=0}^ny_j\ell_j(x),\qquad
 \ell_j(x)=\prod_{i\ne j}\frac{x-x_i}{x_j-x_i}.
 \label{eq:na-lagrange}
\end{equation}
\end{theorem}
\begin{proof}
$\ell_j(x_i)$ 在 $i=j$ 时为 $1$，其余为 $0$，故式~\eqref{eq:na-lagrange} 满足所有条件。若有两解，其差为次数不超过 $n$ 且具有 $n+1$ 个不同零点的多项式，因此恒为零。
\end{proof}

在单项式基下直接解 Vandermonde 系统也可得到插值多项式，但该表示可能很病态。\keyterm{重心插值公式（barycentric interpolation formula）}令 $w_j=1/\prod_{i\ne j}(x_j-x_i)$，则当 $x$ 不等于任一节点时
\begin{equation}
 p_n(x)=\frac{\sum_{j=0}^n w_jy_j/(x-x_j)}{\sum_{j=0}^n w_j/(x-x_j)}.
 \label{eq:na-barycentric}
\end{equation}
推导只需把式~\eqref{eq:na-lagrange} 中共同因子 $\prod_i(x-x_i)$ 约去，并利用 $\sum_j\ell_j(x)=1$。在节点处直接返回对应数据，避免除零。权重可共同乘任意非零常数，预处理后每次评价为 $O(n)$；一般节点的朴素权重计算为 $O(n^2)$，大量节点还需要防止权重溢出或下溢。

\section{差商、Newton 形式与余项}
差商（divided difference）递归定义为 $f[x_i]=f(x_i)$ 以及
\[
 f[x_i,\ldots,x_{i+k}]=
 \frac{f[x_{i+1},\ldots,x_{i+k}]-f[x_i,\ldots,x_{i+k-1}]}{x_{i+k}-x_i}.
\]
插值多项式可写成
\[
 p_n(x)=f[x_0]+\sum_{k=1}^nf[x_0,\ldots,x_k]\prod_{j=0}^{k-1}(x-x_j).
\]
此形式称 Newton 插值形式，增加一个节点时只需增加一个最高次项；嵌套乘法可在 $O(n)$ 时间评价。它与 Newton 求根法是不同问题中的不同算法。

\begin{theorem}[插值余项]
\label{thm:na-interpolation-error}
设 $f\in C^{n+1}[a,b]$，节点均属于 $[a,b]$。对任意 $x\in[a,b]$，有某个 $\xi$ 位于节点与 $x$ 的最小覆盖区间内，使
\begin{equation}
 f(x)-p_n(x)=\frac{f^{(n+1)}(\xi)}{(n+1)!}\prod_{j=0}^n(x-x_j).
 \label{eq:na-interpolation-error}
\end{equation}
在节点处等式两边均为零，不必指定唯一的 $\xi$。
\end{theorem}
\begin{proof}
对非节点 $x$，记 $\omega(t)=\prod_j(t-x_j)$，取 $C=(f(x)-p_n(x))/\omega(x)$。函数 $g(t)=f(t)-p_n(t)-C\omega(t)$ 在所有节点及 $x$ 上为零，连续应用 Rolle 定理得到某点 $g^{(n+1)}(\xi)=0$。由于 $p_n^{(n+1)}=0$、$\omega^{(n+1)}=(n+1)!$，即得余项。
\end{proof}

\begin{example}[二次函数的线性插值]
用节点 $0,1$ 插值 $f(x)=x^2$ 得 $p_1(x)=x$。误差为 $x(x-1)$，在 $[0,1]$ 上最大绝对值为 $1/4$。余项公式也给出 $f''/2!=1$，因而在本例是精确误差表达式。
\end{example}

\section{节点选择、Runge 现象与 Chebyshev 多项式}
\label{sec:na-chebyshev}
当次数增长时，余项同时涉及高阶导数和节点乘积，不能仅根据分母中的阶乘断言误差趋零。函数 $f(x)=1/(1+25x^2)$ 在 $[-1,1]$ 上的等距高次插值会在端部出现严重振荡，这是经典的\keyterm{Runge 现象（Runge phenomenon）}。本书将其用作可复现实验；改善节点分布或采用分段低次方法比盲目加大全局次数更可靠。

\begin{definition}[Chebyshev 多项式]
令 $T_0(x)=1$、$T_1(x)=x$，并递推 $T_{n+1}(x)=2xT_n(x)-T_{n-1}(x)$。由余弦加法公式，$T_n(\cos\theta)=\cos(n\theta)$。$n\ge1$ 时其最高次项系数为 $2^{n-1}$，且在 $[-1,1]$ 上绝对值不超过 $1$。
\end{definition}

\begin{proposition}[节点乘积的极小性质]
在所有 $n+1$ 次首一多项式中，$2^{-n}T_{n+1}$ 在 $[-1,1]$ 上达到最小的一致范数 $2^{-n}$。因而选择节点
\[
 x_j=\cos\frac{(2j+1)\pi}{2(n+1)},\qquad j=0,\ldots,n,
\]
可使节点乘积的最大绝对值最小。
\end{proposition}
\begin{proof}
该首一多项式在 $n+2$ 个余弦极值点交替取 $\pm2^{-n}$。若另一个首一多项式在整个区间绝对值严格小于 $2^{-n}$，二者之差会在相邻极值点之间至少有一个零点，共至少 $n+1$ 个；但最高次项相消后其次数至多为 $n$，矛盾。
\end{proof}

节点的误差放大还可用 Lebesgue 函数 $\Lambda_n(x)=\sum_j|\ell_j(x)|$ 描述。若数据误差满足 $|\Delta y_j|\le\delta$，则插值值误差至多为 $\Lambda_n(x)\delta$。令 $\Lambda_n=\max_{[-1,1]}\Lambda_n(x)$，任取 $q\in\mathcal P_n$，由插值算子 $I_nq=q$ 得
\[
 \|f-I_nf\|_\infty\le(1+\Lambda_n)\|f-q\|_\infty.
\]
取所有 $q$ 的下确界，说明插值误差由函数本身的可逼近性和节点造成的放大共同决定。Chebyshev 节点改善一般插值行为，但对任意连续函数都无条件一致收敛的说法仍然过强。

\section{分段线性、Hermite 插值与三次样条}
设 $a=x_0<\cdots<x_N=b$，$h_i=x_{i+1}-x_i$，$h=\max_i h_i$。每个小区间使用两端点线性插值，拼成连续的分段线性函数 $I_hf$。若 $f\in C^2[a,b]$，定理~\ref{thm:na-interpolation-error} 给出
\[
 \|f-I_hf\|_\infty\le\frac{h^2}{8}\|f''\|_\infty,
\]
因为每段上 $|(x-x_i)(x-x_{i+1})|\le h_i^2/4$。局部低次逼近随着网格加密收敛，不需要不断提高单段次数。

若端点导数也已知，可以在一段 $[a,b]$ 上构造三次 Hermite 插值。令 $t=(x-a)/(b-a)$，则
\begin{align*}
 H(x)={}&(2t^3-3t^2+1)f(a)+(-2t^3+3t^2)f(b)\\
 &+(b-a)(t^3-2t^2+t)f'(a)+(b-a)(t^3-t^2)f'(b).
\end{align*}
四个端点条件直接确定唯一三次多项式。若 $f\in C^4[a,b]$，重复零点版本的 Rolle 论证给出误差 $f(x)-H(x)=f^{(4)}(\xi)(x-a)^2(x-b)^2/4!$。这里导数数据若由差分估计而来，其误差必须另计。

\keyterm{三次样条（cubic spline）}在各小区间为三次以下多项式，并要求整体二阶连续。只指定节点值还不足以唯一确定样条，需要两个边界条件。自然边界条件为 $S''(a)=S''(b)=0$；夹持边界条件给定 $S'(a)$、$S'(b)$，二者含义不同。

设 $M_i=S''(x_i)$。自然三次样条令 $M_0=M_N=0$，其余未知量满足
\begin{equation}
 h_{i-1}M_{i-1}+2(h_{i-1}+h_i)M_i+h_iM_{i+1}
 =6\left(\frac{y_{i+1}-y_i}{h_i}-\frac{y_i-y_{i-1}}{h_{i-1}}\right),
 \quad 1\le i<N.
 \label{eq:na-spline-system}
\end{equation}
在 $[x_i,x_{i+1}]$ 内，函数由
\begin{align*}
 S(x)={}&\frac{M_i(x_{i+1}-x)^3+M_{i+1}(x-x_i)^3}{6h_i}\\
 &+\left(y_i-\frac{M_ih_i^2}{6}\right)\frac{x_{i+1}-x}{h_i}
 +\left(y_{i+1}-\frac{M_{i+1}h_i^2}{6}\right)\frac{x-x_i}{h_i}
\end{align*}
给出。它的二阶导数是端点二阶导数的线性插值，积分两次并代入节点值即可推导；再令相邻段的一阶导数连续便得到式~\eqref{eq:na-spline-system}。

内部线性系统对称且严格对角占优，对角元为正，故正定并唯一可解，可用三对角方法以 $O(N)$ 代价求出全部 $M_i$。样条的整体误差与函数光滑性、边界条件及网格有关；自然边界若与原函数端点二阶导数不一致，不能直接套用匹配边界条件时的全区间四阶误差断言。

\section{习题}
\begin{enumerate}[label=\arabic*.]
\item\label{exr:na-interp-poly} 用节点 $0,1,2$ 与数据 $1,2,5$，分别写出 Lagrange 和 Newton 插值形式，并求 $p(3/2)$。
\item\label{exr:na-interp-error} 对 $f(x)=\exp x$ 在 $[0,1]$ 的均匀分段线性插值，给出保证一致误差不超过 $10^{-4}$ 的网格段数充分条件。
\item\label{exr:na-interp-spline} 求节点 $0,1,2$、数据 $0,1,0$ 对应的自然三次样条的节点二阶导数，并写出第一段表达式。
\item\label{exr:na-interp-cheb} 求 $T_3$，列出它的零点，核对首一化后的最大绝对值。
\end{enumerate}

\chapter{数值微分与数值积分}
\label{chap:na-quadrature}

\section{差分公式与误差平衡}
\keyterm{有限差分（finite difference）}用函数值的线性组合近似导数。Taylor 展开给出前向差分 $D_h^+f(x)=[f(x+h)-f(x)]/h=f'(x)+O(h)$，以及中心差分
\begin{equation}
 D_h^0f(x)=\frac{f(x+h)-f(x-h)}{2h}=f'(x)+O(h^2).
 \label{eq:na-central-difference}
\end{equation}
前者要求局部二阶光滑，后者要求局部三阶光滑。若 $f\in C^3[x-h,x+h]$，中心差分的截断误差绝对值不超过 $h^2\|f'''\|_\infty/6$，由两端 Taylor 余项相加可得。

如果每个函数值带有绝对误差至多 $\eta$，中心差分的输入误差贡献至多为 $\eta/h$。因此差分越细并非越准：截断误差随 $h$ 减小，数据与舍入误差却被除以 $h$ 放大。若 $\eta$ 与 $u$ 同阶且函数尺度固定，平衡 $C_1h^2+C_2u/h$ 得到 $h$ 的量级为 $u^{1/3}$；这只是尺度模型，不是对所有函数通用的最佳步长。

二阶中心差分为
\[
 \frac{f(x+h)-2f(x)+f(x-h)}{h^2}=f''(x)+O(h^2),
\]
通常要求 $f\in C^4$。它把函数值噪声进一步放大到 $O(\eta/h^2)$，因而对含噪测量直接反复数值微分尤其需要谨慎。

\section{求积公式与代数精度}
\begin{definition}[求积公式的代数精度]
对积分 $I(f)=\int_a^bf(x)\,\mathrm dx$，求积公式（quadrature rule）为 $Q(f)=\sum_{j=1}^m w_jf(x_j)$。若对所有次数不超过 $d$ 的多项式有 $Q(p)=I(p)$，但对某个 $d+1$ 次多项式不成立，则称其代数精度（degree of exactness）为 $d$。
\end{definition}
插值型求积先插值再积分，权重为 Lagrange 基函数的积分。梯形公式与 Simpson 公式分别为
\[
 T(f)=\frac{b-a}{2}[f(a)+f(b)],\qquad
 S(f)=\frac{b-a}{6}\left[f(a)+4f\left(\frac{a+b}{2}\right)+f(b)\right].
\]
梯形代数精度为 $1$。Simpson 虽来自二次插值，但对称性使三次项也积得精确，其代数精度为 $3$。代数精度不是对任意函数都精确的声明，也不同于网格方法的全局收敛阶。

\begin{theorem}[复合梯形与复合 Simpson 的误差]
\label{thm:na-composite-quadrature}
令 $x_i=a+ih$、$h=(b-a)/N$。复合梯形
\[
 T_h=h\left[\frac{f(a)+f(b)}2+\sum_{i=1}^{N-1}f(x_i)\right]
\]
在 $f\in C^2[a,b]$ 时满足
\[
 I(f)-T_h=-\frac{b-a}{12}h^2f''(\xi)
\]
对某个 $\xi\in(a,b)$ 成立。若 $N$ 为偶数且 $f\in C^4[a,b]$，复合 Simpson
\[
 S_h=\frac h3\left[f(a)+f(b)+4\sum_{\substack{1\le i<N\\i\text{ 为奇数}}}f(x_i)
 +2\sum_{\substack{2\le i<N\\i\text{ 为偶数}}}f(x_i)\right]
\]
满足 $I(f)-S_h=-(b-a)h^4f^{(4)}(\xi)/180$，其中 $\xi$ 可与前式不同。
\end{theorem}
\begin{proof}[误差公式的推导要点]
在单个宽为 $h$ 的区间，积分线性插值余项并利用节点乘积恒为非正，可得梯形误差 $-h^3f''(\xi_i)/12$。相加后由 $f''$ 连续性得到全局公式。Simpson 误差需要保留其三次精确性，不能仅积分二次插值余项便宣称四阶。令单段误差泛函为 $L=I-S$；由带积分余项的三阶 Taylor 公式及 $L$ 消去三次以下多项式，得 $L(f)=\int_a^b K(t)f^{(4)}(t)\,\mathrm dt$，其中 $K(t)=L((x-t)_+^3/6)$，$(s)_+=\max(s,0)$。在单位区间上直接代入有
\[
 K_0(t)=\begin{cases}
 -t^3(2-3t)/72,&0\le t\le1/2,\\
 -(1-t)^3(3t-1)/72,&1/2\le t\le1.
 \end{cases}
\]
所以 $K_0\le0$，积分为 $-1/2880$。仿射变换给出一般区间上的 $K(t)=(b-a)^4K_0((t-a)/(b-a))$，故 $\int_a^bK(t)\,\mathrm dt=-(b-a)^5/2880$。因此每个宽 $2h$ 的 Simpson 段误差为 $-h^5f^{(4)}(\xi_i)/90$，相加即得全局公式。
\end{proof}

\begin{example}[带确定误差界的积分]
对 $\int_0^1e^x\,\mathrm dx=e-1$，取 $N=10$，复合梯形误差上界为 $e/1200$，复合 Simpson 误差上界为 $e/(180\cdot10^4)$。二者都以同一最细网格为基础；这是可计算的先验界，实际误差通常更小，但不能仅用“通常”代替精度保证。
\end{example}

\section{Richardson 外推、Romberg 积分与自适应}
\label{sec:na-richardson}
若近似量满足 $A(h)=A+ch^p+O(h^{p+q})$，其中 $q>0$，则
\begin{equation}
 \frac{2^pA(h/2)-A(h)}{2^p-1}=A+O(h^{p+q}).
 \label{eq:na-richardson}
\end{equation}
这是消去首个误差项的 Richardson 外推（Richardson extrapolation）。它依赖可用的渐近展开；函数不光滑、含奇点或舍入主导时，形式消项不代表真实改进。

对足够光滑的积分函数，复合梯形误差具有偶次幂展开。令 $R_{j,0}=T_{(b-a)/2^j}$，递推
\[
 R_{j,k}=R_{j,k-1}+\frac{R_{j,k-1}-R_{j-1,k-1}}{4^k-1},\qquad1\le k\le j,
\]
得到 Romberg 表。每一列需要相应更高阶的光滑性才能解释其阶数；复用旧节点，只新增中点评价可以降低函数调用成本。

自适应 Simpson 在一段上比较一次 Simpson 值 $S_1$ 与二等分后的和 $S_2$。若局部四阶误差展开有效，$|S_2-S_1|/15$ 近似估计较细结果的误差，修正值为 $S_2+(S_2-S_1)/15$。算法输入函数、区间、绝对容差和最大递归深度；若估计量未达局部容差，就把容差分给两个子区间并继续，最后相加。必须限制最小区间、递归深度和评价次数。因估计量可能漏掉节点间的尖峰或振荡，它是后验估计而非无条件严格上界。

\section{Gauss 求积：把节点也作为未知量}
令 $w$ 在 $[a,b]$ 上连续且严格为正，考虑 $I_w(f)=\int_a^bf(x)w(x)\,\mathrm dx$。对该权构造 $n$ 次首一正交多项式 $\pi_n$，即它与所有 $\mathcal P_{n-1}$ 中多项式正交。其存在唯一性可由对单项式作正交化得到。

\begin{lemma}[正交多项式的零点]
$\pi_n$ 在 $(a,b)$ 中有 $n$ 个互异实根。
\end{lemma}
\begin{proof}
设区间内部实际发生变号的不同零点为 $\xi_1,\ldots,\xi_k$，令 $q(x)=\prod_{j=1}^k(x-\xi_j)$。若 $k<n$，$\pi_nq$ 在 $(a,b)$ 不变号且不恒为零，因此其正权积分非零；但 $\deg q<n$，正交性又使该积分为零，矛盾。故 $k=n$，次数限制说明所有根均简单且在内部。
\end{proof}

\begin{theorem}[Gauss 公式的精确性与误差]
\label{thm:na-gauss}
取 $\pi_n$ 的根为节点，以插值积分定义权重 $w_j=\int_a^b\ell_j(x)w(x)\,\mathrm dx$，则所得公式对所有次数不超过 $2n-1$ 的多项式精确，权重全部为正。若 $f\in C^{2n}[a,b]$，则
\[
 I_w(f)-\sum_{j=1}^nw_jf(x_j)
 =\frac{f^{(2n)}(\xi)}{(2n)!}\int_a^b\pi_n(x)^2w(x)\,\mathrm dx
\]
对某个 $\xi\in(a,b)$ 成立。
\end{theorem}
\begin{proof}
任取 $p\in\mathcal P_{2n-1}$，作带余除法 $p=q\pi_n+r$，其中 $q,r\in\mathcal P_{n-1}$。正交性使 $q\pi_n$ 的积分为零，节点求和也为零；$r$ 则由插值精确积分。对 $\ell_j^2$ 使用已证精确性，得 $w_j=\int_a^b\ell_j^2w>0$。令 $H$ 为同时匹配各节点函数值和一阶导数的 $2n-1$ 次 Hermite 插值，重复 Rolle 定理给出 $f-H=f^{(2n)}(\xi_x)\pi_n^2/(2n)!$。该余项被 $f^{(2n)}$ 的最小值和最大值乘非负权夹住，积分后由连续性得到一个共同的 $\xi$。因公式对 $H$ 精确，余项就是求积误差。
\end{proof}

代数精度恰为 $2n-1$：对 $\pi_n^2$，节点求和为零，而正权积分严格为正。对 $[-1,1]$ 上权 $1$，两点 Gauss--Legendre 公式为
\[
 \int_{-1}^1f(x)\,\mathrm dx\approx f(-1/\sqrt3)+f(1/\sqrt3).
\]
一般区间用仿射代换 $x=(a+b)/2+(b-a)t/2$，权重也乘 $(b-a)/2$。端点奇点或无限区间要求变量变换、特殊权函数或专用算法，不能直接继承有限闭区间上的光滑误差阶。

\section{习题}
\begin{enumerate}[label=\arabic*.]
\item\label{exr:na-quad-difference} 用 Taylor 展开推导二阶中心差分，并给出 $f\in C^4$ 时其截断误差上界。
\item\label{exr:na-quad-simpson} 在 $[0,1]$ 上对 $x^3$ 和 $x^4$ 使用单段 Simpson，核对精确性和 $I-S$ 的符号。
\item\label{exr:na-quad-gauss} 推导 $[-1,1]$ 上对称两点公式的节点与权重，使其对三次以下多项式精确。
\item\label{exr:na-quad-richardson} 已知 $A(h)=A+ch^2+dh^4+O(h^6)$，写出一次外推并求剩余的 $h^4$ 项系数。
\end{enumerate}

\chapter{常微分方程初值问题}
\label{chap:na-ode}

\section{初值问题与误差传播}
考虑向量初值问题 $y'(t)=F(t,y(t))$、$y(t_0)=y_0\in\mathbb R^d$，目标是在有限区间 $[t_0,T]$ 上近似解。若 $F$ 关于 $t$ 连续、关于 $y$ 在相关区域局部 Lipschitz，则初值附近有唯一局部解；是否延拓到 $T$ 还需排除有限时间发散。本章的收敛讨论均假定所研究的精确解在整个区间存在，且所需光滑性和 Lipschitz 界在包含解与数值轨迹的区域成立。向量与矩阵范数采用相容的选定范数。

若两个精确解 $y,z$ 的初始数据不同，而 $F$ 对 $y$ 的 Lipschitz 常数为 $L$，则积分方程和 Gronwall 不等式给出 $\|y(t)-z(t)\|\le e^{L(t-t_0)}\|y(t_0)-z(t_0)\|$。这说明初值扰动随时间积累，长时间计算的困难可能来自方程本身。数值方法还会在每一步不断注入局部误差。

\section{Euler 方法与全局一阶收敛}
在均匀网格 $t_n=t_0+nh$ 上，显式 Euler 方法为
\begin{equation}
 y_{n+1}=y_n+hF(t_n,y_n).
 \label{eq:na-euler}
\end{equation}
它把一个时间段内的导数冻结在左端点。输入为 $F$、初值、步长和终点，输出为网格近似值；每步一次函数评价，若 $d$ 固定且评价代价固定，总成本与步数同阶。非整步到达终点时应缩短最后一步，而不是越过终点后仍报告为终点解。

\begin{definition}[局部缺陷与全局误差]
对一步法 $y_{n+1}=\Phi_h(t_n,y_n)$，用精确解代入一步得到的局部缺陷为 $d_{n+1}=y(t_{n+1})-\Phi_h(t_n,y(t_n))$；全局误差为 $e_n=y_n-y(t_n)$。本书不把局部缺陷除以 $h$；有些文献把 $d_{n+1}/h$ 称为局部截断误差，比较阶数时需先核对约定。
\end{definition}

\begin{theorem}[Euler 的全局误差界]
若 $F$ 关于 $y$ 的 Lipschitz 常数为 $L\ge0$，$\max_{[t_0,T]}\|y''(t)\|\le M$，且 $y_0=y(t_0)$，则显式 Euler 满足
\[
 \|e_n\|\le\frac{Mh}{2L}\bigl(e^{L(t_n-t_0)}-1\bigr)\quad(L>0),
\]
当 $L=0$ 时右端取 $Mh(t_n-t_0)/2$。
\end{theorem}
\begin{proof}
积分形式的 Taylor 余项给出 $\|d_{n+1}\|\le Mh^2/2$。相减并取范数得 $\|e_{n+1}\|\le(1+hL)\|e_n\|+Mh^2/2$。展开几何级数，得到 $Mh[(1+hL)^n-1]/(2L)$，再用 $1+hL\le e^{hL}$。当 $L=0$ 时直接累加。
\end{proof}

更一般地，若一步映射满足 $\|\Phi_h(t,u)-\Phi_h(t,v)\|\le(1+Ch)\|u-v\|$，局部缺陷一致为 $O(h^{p+1})$，初始误差为 $O(h^p)$，相同论证给出固定时间区间上的全局误差 $O(h^p)$。这明确展示了一致性与误差传播控制怎样共同导致收敛。

隐式 Euler 方法为 $y_{n+1}=y_n+hF(t_{n+1},y_{n+1})$。每一步需要解代数方程；$hL<1$ 时其不动点映射为压缩，在全空间全局 Lipschitz 的设置下可保证该步唯一可解。其他结构可能允许更大步长，故不能把这个充分条件误认为必要条件。隐式性带来的稳定优势需付出求解代数系统的代价。

\section{Runge--Kutta 方法}
显式中点法通过一个中间预测改善导数近似：
\[
 k_1=F(t_n,y_n),\quad k_2=F(t_n+h/2,y_n+hk_1/2),\quad y_{n+1}=y_n+hk_2.
\]
若 $F$ 充分光滑，展开 $k_2=F+\tfrac h2(F_t+F_yF)+O(h^2)$，其更新匹配精确解的二阶 Taylor 展开，故局部缺陷为 $O(h^3)$，在上述 Lipschitz 条件下全局二阶。改进 Euler 又称显式梯形法，使用 $k_2=F(t_n+h,y_n+hk_1)$ 及 $y_{n+1}=y_n+h(k_1+k_2)/2$，同样为二阶，但与中点法不是同一个更新式。

经典四阶 Runge--Kutta（RK4）为
\begin{align}
 k_1&=F(t_n,y_n),\notag\\
 k_2&=F(t_n+h/2,y_n+hk_1/2),\notag\\
 k_3&=F(t_n+h/2,y_n+hk_2/2),\notag\\
 k_4&=F(t_n+h,y_n+hk_3),\notag\\
 y_{n+1}&=y_n+\frac h6(k_1+2k_2+2k_3+k_4).
 \label{eq:na-rk4}
\end{align}
在 $F$ 具有所需连续至四阶偏导且相关导数有界的区域中，展开各阶段并匹配精确解至四次项，可得局部缺陷 $O(h^5)$，因此固定时间区间上全局四阶。这里只说明阶数证明的匹配方法，不列出全部多元 Taylor 项；把该公式直接视为对时间积分使用 Simpson 而忽略阶段值误差，不构成四阶证明。

一般显式 $s$ 阶段方法写为 $k_i=F(t_n+c_ih,y_n+h\sum_{j<i}a_{ij}k_j)$、$y_{n+1}=y_n+h\sum_i b_ik_i$。本书采用 $c_i=\sum_{j<i}a_{ij}$。一阶条件为 $\sum b_i=1$，二阶还需 $\sum b_ic_i=1/2$；高阶需要更多独立条件。方法名称中的“阶段数”和“精度阶数”不是同一概念。

\section{绝对稳定性与刚性}
\label{sec:na-ode-stability}
\begin{definition}[稳定函数与绝对稳定域]
将一步法用于标量测试方程 $y'=\lambda y$、$\lambda\in\mathbb C$，若得到 $y_{n+1}=R(z)y_n$、$z=h\lambda$，则 $R$ 为稳定函数。集合 $\{z:|R(z)|\le1\}$ 称闭绝对稳定域；其严格内部 $|R(z)|<1$ 对应随步数衰减。若该闭域包含整个左半平面 $\operatorname{Re}z\le0$，称方法 A 稳定。
\end{definition}

显式 Euler 有 $R(z)=1+z$，因此对实 $\lambda<0$ 要求 $0<h\le2/|\lambda|$ 才不增长，严格衰减要求严格小于上限。隐式 Euler 有 $R(z)=1/(1-z)$，对 $\operatorname{Re}z\le0$ 有 $|1-z|\ge1$，故 A 稳定。隐式梯形法
\[
 y_{n+1}=y_n+\tfrac h2\bigl[F(t_n,y_n)+F(t_{n+1},y_{n+1})\bigr]
\]
为二阶，其 $R(z)=(1+z/2)/(1-z/2)$ 也 A 稳定，但当实 $z\to-\infty$ 时趋于 $-1$。隐式 Euler 的稳定函数则趋于零，因而更有效地抑制快速衰减模态。A 稳定且在左半平面趋于无穷远时 $R(z)\to0$ 的方法称 L 稳定；隐式 Euler 满足这一性质。

\keyterm{刚性（stiffness）}描述这样一类计算现象：解的主要变化尺度允许较大步长，但显式方法为了稳定必须采用很小步长。它依赖方程、区间、方法和精度要求，不仅由某个系数“大”来定义。

\begin{example}[平滑解与快速扰动模态]
令 $y'=-1000(y-\cos t)-\sin t$，$y(0)=1$，则精确解为 $y(t)=\cos t$。任意偏离该解的扰动满足 $e'=-1000e$。显式 Euler 的稳定限制为 $h\le0.002$，尽管余弦本身在单位时间尺度平缓变化。隐式方法能够缓解稳定步长约束，但步长仍须满足精度要求。
\end{example}

标量稳定域分析是理解线性系统的起点。对可正交对角化的系统，它逐模态适用；对非正规矩阵，特征值位置还不能完全描述短时范数增长。非线性系统则需要局部线性化及更具体的稳定分析。

\section{线性多步法与步长控制}
线性 $s$ 步法写为 $\sum_{j=0}^s\alpha_jy_{n+j}=h\sum_{j=0}^s\beta_jF(t_{n+j},y_{n+j})$，其中 $\alpha_s\ne0$。它复用历史函数值，例如二步 Adams--Bashforth 法为
\[
 y_{n+2}=y_{n+1}+\frac h2(3F_{n+1}-F_n),\qquad F_j=F(t_j,y_j).
\]
初始的 $y_1$ 需要另用足够精度的一步法产生。$\beta_s=0$ 时为显式，否则通常为隐式。

对固定步长与固定系数，一致性至少要求 $\sum\alpha_j=0$、$\sum j\alpha_j=\sum\beta_j$。\keyterm{零稳定性（zero-stability）}要求有限时间内初值及局部扰动的影响在 $h\to0$ 时不会无限放大；对这类线性多步法等价于特征多项式 $\rho(\zeta)=\sum\alpha_j\zeta^j$ 的全部根满足 $|\zeta|\le1$，且单位圆上的根均简单。对满足标准 Lipschitz 条件的初值问题，一致性与零稳定性共同保证收敛，前提是起步值也收敛。这是线性多步法的基本等价定理，本书陈述其使用条件，不展开差分方程理论中的完整证明。

零稳定性研究步长趋零时的误差传播，绝对稳定性研究固定 $h\lambda$ 下的长步数行为，两者不能互相替代。Adams--Bashforth 二步法的 $\rho(\zeta)=\zeta^2-\zeta$，根为 $0,1$，故零稳定；其显式稳定域仍有限。

自适应步长通常比较一对嵌入式方法，形成局部误差估计 $E$。若 $E$ 的主项与 $h^{q+1}$ 成比例，可用安全因子 $0<s<1$ 建议 $h_{\mathrm{new}}=sh(\mathrm{tol}/E)^{1/(q+1)}$，并限制增长与缩小倍数。向量分量宜按 $\varepsilon_{\mathrm{abs},i}+\varepsilon_{\mathrm{rel}}|y_i|$ 归一化。局部容差控制不自动等于最终全局误差保证，尤其在不稳定动力系统或长时间积分中。

\section{习题}
\begin{enumerate}[label=\arabic*.]
\item\label{exr:na-ode-euler} 对 $y'=y$、$y(0)=1$，写出显式与隐式 Euler 在 $t=1$、$h=1/N$ 下的数值值，并判断偏大还是偏小，隐式情形取 $N>1$。
\item\label{exr:na-ode-stability} 对 $y'=-20y$，求显式 Euler 不增长的步长范围，并解释边界步长处的行为。
\item\label{exr:na-ode-rk} 将 RK4 用于 $y'=\lambda y$，推导稳定函数。
\item\label{exr:na-ode-root} 考察多步格式 $y_{n+2}-3y_{n+1}+2y_n=-hF_n$：检查一致性与零稳定性，说明即使局部缺陷趋零也可能不收敛。
\end{enumerate}

\chapter{边值问题与偏微分方程的差分入门}
\label{chap:na-pde}

\section{一维 Poisson 边值问题的离散化}
考虑
\begin{equation}
 -u''(x)=f(x),\quad 0<x<1,\qquad u(0)=\alpha,\quad u(1)=\beta.
 \label{eq:na-poisson}
\end{equation}
当 $f$ 连续时，积分两次并用边界条件确定积分常数，可得唯一二阶连续解。这里的未知是整个区间上的函数，边界条件给在两端，与初值问题不同。

令整数 $N\ge2$、$h=1/N$、$x_i=ih$，以 $U_i$ 近似 $u(x_i)$。中心差分得到
\begin{equation}
 -\frac{U_{i-1}-2U_i+U_{i+1}}{h^2}=f(x_i),\quad1\le i<N,
 \qquad U_0=\alpha,\quad U_N=\beta.
 \label{eq:na-poisson-discrete}
\end{equation}
内部未知量组成 $N-1$ 维向量，系数矩阵为 $A_h=h^{-2}\operatorname{tridiag}(-1,2,-1)$，右端首、末分量分别加入 $\alpha/h^2$、$\beta/h^2$；只有一个内部点时，两项加到同一分量。对补零边界的任意内部向量 $v$，
\[
 v^{\mathsf T}A_hv=h^{-2}\sum_{i=0}^{N-1}(v_{i+1}-v_i)^2>0\quad(v\ne0),
\]
故 $A_h$ 为 SPD，离散系统唯一可解。

\section{离散最大值原理与二阶误差}
\begin{lemma}[离散比较原理]
\label{lem:na-discrete-maximum}
若网格函数 $v$ 满足 $v_0,v_N\ge0$ 且 $2v_i-v_{i-1}-v_{i+1}\ge0$，则所有 $v_i\ge0$。
\end{lemma}
\begin{proof}
若存在负最小值，它只能在内部达到。在最小点，两侧值不小于该值，而给定不等式迫使两侧都等于它。递推到边界仍为同一负值，与边界非负矛盾。
\end{proof}

\begin{theorem}[Poisson 差分格式的全局误差]
若式~\eqref{eq:na-poisson} 的解 $u\in C^4[0,1]$，则式~\eqref{eq:na-poisson-discrete} 的解满足
\[
 \max_{0\le i\le N}|U_i-u(x_i)|\le\frac{h^2}{96}\|u^{(4)}\|_\infty.
\]
\end{theorem}
\begin{proof}
以 $u_h$ 表示精确解的内部采样，差分 Taylor 余项给出缺陷 $\tau=A_hu_h-b_h$ 满足 $\|\tau\|_\infty\le h^2\|u^{(4)}\|_\infty/12$，其中 $b_h$ 已包含边界贡献。误差 $e=U-u_h$ 满足 $A_he=-\tau$ 且端点为零。屏障函数 $w_i=x_i(1-x_i)/2$ 满足 $A_hw=1$、$w_0=w_N=0$。对 $\|\tau\|_\infty w\pm e$ 应用引理~\ref{lem:na-discrete-maximum}，得到 $|e_i|\le\|\tau\|_\infty w_i\le\|\tau\|_\infty/8$。
\end{proof}

这个证明分为缺陷估计和离散求解算子的稳定控制两步。仅把精确解代入格式发现残差为 $O(h^2)$，还不足以证明解误差为 $O(h^2)$；必须控制残差到解的放大。

\begin{example}[制造解检验]
选取 $u(x)=\sin(\pi x)$，则 $f(x)=\pi^2\sin(\pi x)$、边界为零。把这一人为选择的精确解对应的数据代入程序，逐次减半 $h$，在离散误差主导时最大误差之比应趋近 $4$。这称为制造解法（method of manufactured solutions），用于检验实现与预期阶数，并不表示所有真实问题都具有已知精确解。
\end{example}

\section{热方程与时间、空间步长的耦合}
对 $u_t=\nu u_{xx}$、$\nu>0$，在 $[0,1]$ 上配合已给定的初值与 Dirichlet 边界，采用空间步长 $h$、时间步长 $\Delta t$，得到显式格式
\begin{equation}
 U_i^{n+1}=rU_{i-1}^n+(1-2r)U_i^n+rU_{i+1}^n,
 \qquad r=\frac{\nu\Delta t}{h^2}.
 \label{eq:na-heat-explicit}
\end{equation}
当 $0\le r\le1/2$ 时，各系数非负且和为 $1$，故无源项、齐次边界下最大绝对值不会增加。周期网格上的模态 $e^{\mathrm i i\theta}$ 有放大因子 $G(\theta)=1-4r\sin^2(\theta/2)$，要求所有频率不增长也得到 $r\le1/2$。有限 Dirichlet 网格可有略宽的精确谱阈值，但 $1/2$ 是与网格无关的标准充分条件。

若精确解具有有界的二阶时间导数和四阶空间导数，按 $\Delta t$ 归一化的局部缺陷为 $O(\Delta t+h^2)$。在上述最大范数稳定条件下，把每一步实际注入的 $\Delta t\,O(\Delta t+h^2)$ 累加，在固定终止时间可得全局误差 $O(\Delta t+h^2)$，还需初始和边界离散误差不超过同阶。

\begin{example}[稳定约束的实际代价]
取 $\nu=1$、$h=1/100$，显式格式要求 $\Delta t\le1/20000$。即使解变化平缓，计算到 $t=1$ 仍至少需要 $20000$ 步。将空间步长减半，允许时间步长变成原来的四分之一；在一维中网格点也约加倍，总工作量可能约增为八倍。
\end{example}

\section{隐式格式、方法线与适用范围}
对齐次边界，半离散系统为 $U'(t)=-\nu A_hU(t)$。先离散空间、再用 ODE 方法积分的策略称为\keyterm{方法线（method of lines）}。随 $h\to0$，$A_h$ 最大特征值与 $h^{-2}$ 同阶，因此系统出现快速衰减尺度。

隐式 Euler 给出 $(I+\nu\Delta t A_h)U^{n+1}=U^n$，每步是 SPD 三对角系统。它不受显式格式的同一稳定步长限制，但时间仍为一阶。Crank--Nicolson 格式为
\[
 (I+\tfrac12\nu\Delta t A_h)U^{n+1}
 =(I-\tfrac12\nu\Delta t A_h)U^n.
\]
在解足够光滑且边界处理一致时，它在时间和空间均为二阶。由于对应隐式梯形法，高频模态在过大时间步下可能产生衰减很慢的符号振荡，二范数稳定并不等于保持非负性或单调性。

本章只处理一维二阶椭圆问题和线性热方程的基础差分格式。一般区域、对流主导问题、复杂边界及非线性守恒律需要有限元、有限体积、上风方法或其他专门理论，不能把此处的稳定与收敛条件直接照搬。

\section{习题}
\begin{enumerate}[label=\arabic*.]
\item\label{exr:na-pde-exact} 对 $-u''=2$、$u(0)=u(1)=0$，求精确解并证明中心差分在节点上精确恢复它。
\item\label{exr:na-pde-spectrum} 验证 $A_h$ 的特征向量分量为 $\sin(k\pi i/N)$，求相应特征值，$1\le k<N$。
\item\label{exr:na-pde-stability} 在周期偶数网格上取最高频误差 $U_i^0=(-1)^i$，分析 $r=0.6$ 时的显式热方程迭代。
\item\label{exr:na-pde-mesh} 为使显式热方程格式的时间和空间误差均与 $h^2$ 同阶，应如何选择 $\Delta t$？还需检查什么约束？
\end{enumerate}

\chapter{离散傅里叶变换与快速算法}
\label{chap:na-fourier}

\section{周期采样与离散正交性}
周期现象适合用三角函数而非全局单项式表示。本章把长度为 $N$ 的复数数据 $v_0,\ldots,v_{N-1}$ 理解为周期区间 $[0,2\pi)$ 上的采样，节点 $x_j=2\pi j/N$。规定离散傅里叶变换（discrete Fourier transform，DFT）的归一化为
\begin{equation}
 \widehat v_k=\sum_{j=0}^{N-1}v_j e^{-2\pi\mathrm i jk/N},\quad
 v_j=\frac1N\sum_{k=0}^{N-1}\widehat v_k e^{2\pi\mathrm i jk/N}.
 \label{eq:na-dft}
\end{equation}
不同软件的归一化可能不同，应先核对正变换与逆变换的系数和指数符号。

\begin{proposition}[离散正交与 Parseval 等式]
对整数 $k,\ell$，有
\[
 \sum_{j=0}^{N-1}e^{2\pi\mathrm i j(k-\ell)/N}
 =\begin{cases}N,&k\equiv\ell\pmod N,\\0,&\text{其余情形}.
 \end{cases}
\]
因此式~\eqref{eq:na-dft} 的逆公式成立，且 $\sum_j|v_j|^2=N^{-1}\sum_k|\widehat v_k|^2$。
\end{proposition}
\begin{proof}
当公比为 $1$ 时求和为 $N$，否则用几何级数公式，分子为 $1-e^{2\pi\mathrm i(k-\ell)}=0$，分母非零。把正变换代入逆变换可恢复每个 $v_j$。再把逆变换代入平方和，并利用共轭和上述正交关系，交叉项全部消失，得到 Parseval 等式。
\end{proof}

\section{混叠与频率解释}
频率 $k$ 与 $k+qN$ 在网格上取值相同，这称为\keyterm{混叠（aliasing）}。仅由这组样本无法区分它们；增加输出曲线上的绘制点不能补回已丢失的采样信息。实数据满足 $\widehat v_{N-k}=\overline{\widehat v_k}$，索引按模 $N$ 解释，因而后一半索引通常代表负频率。

为避免偶数点数的最高频模式约定问题，先设 $N=2m+1$ 为奇数，用频率 $-m,\ldots,m$，令 $c_k=\widehat v_{k\bmod N}/N$。三角插值多项式为 $p(x)=\sum_{k=-m}^{m}c_ke^{\mathrm i kx}$，在全部节点匹配数据。偶数 $N$ 时的 Nyquist 频率 $N/2$ 与 $-N/2$ 在样本上重合，求导等运算需要明确处理这一特殊模态，不能把数组索引全部视为非负物理频率。

\begin{example}[样本无法识别的两种振荡]
在 $N=8$ 的网格上，$e^{\mathrm i x}$ 与 $e^{9\mathrm i x}$ 的取值完全相同，因为 $e^{8\mathrm i x_j}=1$。它们在采样点之外不同，导数也不同。以样本作插值隐含了对可表示频率范围的限制。
\end{example}

\section{快速傅里叶变换与卷积}
直接计算 DFT 需要 $O(N^2)$ 次复运算。若 $N$ 为偶数，把下标分成偶数和奇数，定义长度 $N/2$ 的两个 DFT 为 $E_k,O_k$，则
\[
 \widehat v_k=E_k+e^{-2\pi\mathrm i k/N}O_k,\qquad
 \widehat v_{k+N/2}=E_k-e^{-2\pi\mathrm i k/N}O_k,
 \quad0\le k<N/2.
\]
当 $N=2^s$ 时递归应用，运算量满足 $T(N)=2T(N/2)+O(N)$，故为 $O(N\log N)$。这类计算 DFT 的快速算法称 FFT；FFT 改变计算方式，并没有改变所定义的变换或提高采样所含信息量。

周期卷积定义为 $c_j=\sum_{\ell=0}^{N-1}a_\ell b_{(j-\ell)\bmod N}$。代入 DFT 并重新编号有限求和可得 $\widehat c_k=\widehat a_k\widehat b_k$，故可通过两次正变换、逐点相乘和一次逆变换求卷积。若要求长度 $M$ 与 $K$ 的序列的普通线性卷积，应先补零至至少 $M+K-1$，否则周期绕回会污染结果。

\section{谱微分与周期梯形积分}
对上述奇数节点的三角插值，有 $p'(x)=\sum_{k=-m}^{m}(\mathrm i k)c_ke^{\mathrm i kx}$。因此谱微分可以通过 DFT、频率乘子和逆 DFT 实现。高频乘子 $|k|$ 同时放大高频噪声，所以快速或高阶方法不意味着对含噪数据稳定地求导。

若 $2\pi$ 周期函数有绝对可和的傅里叶级数 $f(x)=\sum_{k\in\mathbb Z}c_ke^{\mathrm i kx}$，则可以交换有限节点求和与级数，得到周期梯形值
\[
 Q_N=\frac{2\pi}{N}\sum_{j=0}^{N-1}f(x_j)=2\pi\sum_{q\in\mathbb Z}c_{qN},\qquad
 \int_0^{2\pi}f(x)\,\mathrm dx=2\pi c_0.
\]
由此 $|Q_N-I|\le2\pi\sum_{q\ne0}|c_{qN}|$。若 $|c_k|\le C|k|^{-s}$、$s>1$，误差为 $O(N^{-s})$；若系数指数衰减，则积分误差也指数衰减。这解释了光滑周期问题中梯形法可能远快于一般二阶估计的原因。端点不匹配的周期延拓会降低光滑性，不能期待同样速度。

\section{习题}
\begin{enumerate}[label=\arabic*.]
\item\label{exr:na-fft-transform} 按式~\eqref{eq:na-dft} 计算序列 $(1,0,-1,0)$ 的 DFT，并用逆变换核验。
\item\label{exr:na-fft-alias} 在 $N=10$ 的网格上，频率 $7$ 与哪个位于 $[-5,4]$ 中的频率不可区分？
\item\label{exr:na-fft-convolution} 计算 $(1,2)$ 与 $(3,4)$ 的线性卷积，并说明直接做长度为 $2$ 的周期卷积得到什么。
\item\label{exr:na-fft-periodic} 对整数 $m\ge1$，求周期梯形公式作用于 $\cos(mx)$ 的结果，并讨论 $N$ 是否整除 $m$。
\end{enumerate}

\chapter{蒙特卡洛积分与随机误差}
\label{chap:na-montecarlo}

\section{把积分表示为期望}
高维积分的规则网格点数随维度快速增长。\keyterm{蒙特卡洛方法（Monte Carlo method）}用随机样本估计期望；其误差是随机变量，必须与确定性截断误差区分。设 $X_1,\ldots,X_N$ 独立同分布于 $[a,b]$ 上的均匀分布，定义
\begin{equation}
 \widehat I_N=\frac{b-a}{N}\sum_{j=1}^Nf(X_j).
 \label{eq:na-mc}
\end{equation}
随机变量用大写表示，一次具体运行得到的 $x_j$ 是它们的取值。伪随机生成器给出可复现的确定性序列，理论分析通常以理想独立随机样本建模；固定种子用于复现实验，不会证明样本独立。

\begin{theorem}[蒙特卡洛估计的均值和方差]
若 $f$ 平方可积，令 $\sigma^2=\operatorname{Var}(f(X_1))$，则 $\mathbb E\widehat I_N=I$ 且
\[
 \operatorname{Var}(\widehat I_N)=\frac{(b-a)^2\sigma^2}{N},\qquad
 \bigl(\mathbb E|\widehat I_N-I|^2\bigr)^{1/2}=\frac{(b-a)\sigma}{\sqrt N}.
\]
若进一步 $\sigma>0$，中心极限定理给出 $\sqrt N(\widehat I_N-I)/[(b-a)\sigma]$ 依分布趋于标准正态。
\end{theorem}
\begin{proof}
均匀分布密度为 $1/(b-a)$，故 $\mathbb E f(X_1)=I/(b-a)$。独立性使不同样本的协方差为零，平均值的方差为单样本方差除以 $N$；估计无偏，均方误差等于方差。最后一句直接使用独立同分布、有限正方差条件下的中心极限定理。
\end{proof}

令 $\overline f=N^{-1}\sum_jf(X_j)$、$s^2=(N-1)^{-1}\sum_j(f(X_j)-\overline f)^2$，其中 $N\ge2$。大样本下可报告近似 $95\%$ 置信区间 $\widehat I_N\pm1.96(b-a)s/\sqrt N$。它是重复抽样意义下的渐近覆盖陈述，既不是本次误差必小于半宽的保证，也不表示未知固定积分值在给定区间内具有由频率学派赋予的后验概率。

\begin{example}[估计平方函数的积分]
对 $I=\int_0^1x^2\,\mathrm dx=1/3$，$\mathbb E X^4=1/5$，所以 $\operatorname{Var}(X^2)=1/5-1/9=4/45$。估计标准差为 $2/\sqrt{45N}$。把标准差缩小到原来的十分之一，需要样本量增加到一百倍。
\end{example}

\section{重要性抽样与方差存在条件}
设 $p$ 是积分区域上的概率密度，在 $f\ne0$ 的地方几乎处处为正。若 $X_j$ 独立服从 $p$，则
\[
 I=\mathbb E_p\frac{f(X)}{p(X)},\qquad
 \widehat I_N^{\mathrm{IS}}=\frac1N\sum_{j=1}^N\frac{f(X_j)}{p(X_j)}.
\]
这里 $\int|f|<\infty$ 保证期望有意义，而有限方差还要求 $\int f(x)^2/p(x)\,\mathrm dx<\infty$。重要性密度在重要区域过小会导致权重爆炸，不能只因抽样“更集中”就断言降低方差。

若 $Z=\int|f|>0$ 已知，理想密度 $p_*=|f|/Z$ 最小化该类估计器的二阶矩：Cauchy--Schwarz 不等式给出 $(\int|f|)^2\le(\int f^2/p)(\int p)=\int f^2/p$，等号条件正是 $p\propto|f|$。若 $f$ 非负，该选择产生零方差，但归一化常数正是待求积分，因此通常只是设计方向，不能当作通用的可执行解法。

\section{控制变量、对偶变量与分层}
若 $Y=f(X)$，另有与它相关且期望 $\mu_Z$ 已知的变量 $Z$，则控制变量估计使用 $Y-c(Z-\mu_Z)$，其期望仍为 $\mathbb EY$。若二者二阶矩有限且 $\operatorname{Var}(Z)>0$，方差关于 $c$ 的二次函数在
\[
 c_* =\frac{\operatorname{Cov}(Y,Z)}{\operatorname{Var}(Z)}
\]
处最小。$c$ 由同一批样本估计时，简单的固定系数无偏证明不再直接适用；可用独立预实验估计参数，再固定它处理正式样本。

对 $U\sim\operatorname{Unif}[0,1]$，用一对 $f(U),f(1-U)$ 的平均构造对偶变量估计。其方差为 $\tfrac12\operatorname{Var}(f(U))+\tfrac12\operatorname{Cov}(f(U),f(1-U))$。在相同两次函数评价成本下，只有协方差为负时才优于两个独立样本的平均；不能仅因为样本成对就保证改善。

将区间划分为互不重叠的子区间，在第 $\ell$ 层独立取 $n_\ell$ 个均匀样本，层长度为 $h_\ell$、函数值方差为 $\sigma_\ell^2$，则分层积分估计的方差为 $\sum_\ell h_\ell^2\sigma_\ell^2/n_\ell$。在连续放松的固定总样本量约束下，最优分配满足 $n_\ell\propto h_\ell\sigma_\ell$；实际需要整数取整并用预实验估计未知方差。

\section{高维与相关样本的边界}
对 $d$ 维区域上的均匀抽样，式~\eqref{eq:na-mc} 中的区间长度替换为体积即可。$N^{-1/2}$ 的指数不显式依赖维数，但方差常数、函数评价成本和生成合适样本的难度都可能随维度恶化。因此“蒙特卡洛没有维数灾难”只能理解为特定误差率中的指数性质，不能理解为高维积分一律容易。

马尔可夫链蒙特卡洛生成相关样本，不能直接套用独立样本的方差证明。对平稳序列 $Y_j$，若自协方差 $\gamma_k=\operatorname{Cov}(Y_1,Y_{1+k})$ 存在，则有限样本恒等式为
\[
 \operatorname{Var}\left(\frac1N\sum_{j=1}^NY_j\right)
 =\frac1N\left[\gamma_0+2\sum_{k=1}^{N-1}\left(1-\frac{k}{N}\right)\gamma_k\right].
\]
由方差双重求和按时间差分组即可推导。进一步使用极限定理还需遍历性及相应依赖条件；平稳或“看起来已混合”本身不自动保证普通中心极限定理。具体采样算法与诊断属于后续概率计算课程，本章保留独立抽样作为可完整分析的基础。

\section{习题}
\begin{enumerate}[label=\arabic*.]
\item\label{exr:na-mc-variance} 用均匀样本估计 $\int_0^1x\,\mathrm dx$，求均值、方差和均方根误差。
\item\label{exr:na-mc-antithetic} 对同一积分使用对偶变量对 $(U,1-U)$，证明每一对得到的估计恰为积分值。
\item\label{exr:na-mc-control} 取 $Y=U^2$、$Z=U$，求最佳控制变量系数以及单样本修正量的方差。
\item\label{exr:na-mc-importance} 对 $f(x)=1$、$0<x<1$，取重要性密度 $p(x)=2x$。证明估计无偏但方差无穷，并说明普通有限方差中心极限定理不能直接使用。
\end{enumerate}

\chapter{科学计算实验与结果验证}
\label{chap:na-practice}

\section{从理论公式到可复现报告}
一个完整计算实验应写明数学问题、输入数据、算法、容差、精度类型与停止原因，再报告误差或残差、运算规模和运行环境。若比较两个方法，函数评价次数、矩阵分解次数或总耗时应采用一致口径。未经测量的理论复杂度不能直接当作实测速度；单次计时也不能代表稳定的性能结论。

误差曲线必须注明纵轴是绝对误差、相对误差、残差还是误差估计量。网格收敛实验应区分离散误差主导区与舍入误差平台；随机实验则需要多次独立重复，而不是只从一个种子的波动轨迹判断收敛率。对于未知真解的问题，可结合更精细计算、守恒关系、独立算法和制造解测试，但这些证据的作用不同，不宜把相互吻合称为无条件证明。

本章代码使用 Python 标准库，作为有限规模、有限实数数据的教学实现。代码中的浮点零、有限性检查和异常信息是程序状态，不替代正文中的实数算术定理。代码示例聚焦算法核心，不承担任意精度计算、区间算术或生产级求解器的全部功能。

\section{包根求解的教学实现}
以下函数接收连续实函数及升序包根区间；函数值的符号需要可靠。返回值依次为中点、当前半区间宽度和停止状态。状态 \texttt{width} 表示通过位置容差停止；\texttt{computed\_zero} 仅表示程序评价得到零。函数并不把它与经过严格区间计算认证的精确零混为一谈。
\begin{verbatim}
import math

def bisect(f, a, b, atol=1e-12, rtol=1e-10, maxiter=200):
    if not (math.isfinite(a) and math.isfinite(b) and a < b):
        raise ValueError("finite ordered endpoints required")
    if not (math.isfinite(atol) and math.isfinite(rtol)
            and atol > 0 and rtol >= 0
            and isinstance(maxiter, int) and maxiter > 0):
        raise ValueError("invalid tolerances or iteration limit")

    def value(x):
        y = f(x)
        if not math.isfinite(y):
            raise ValueError("nonfinite function value")
        return y

    fa, fb = value(a), value(b)
    if fa == 0:
        return a, 0.0, "computed_zero"
    if fb == 0:
        return b, 0.0, "computed_zero"
    if (fa > 0) == (fb > 0):
        raise ValueError("opposite signs required")
    for _ in range(maxiter):
        c = a / 2 + b / 2
        radius = b / 2 - a / 2
        if radius <= atol + rtol * abs(c):
            return c, radius, "width"
        if c == a or c == b:
            raise ArithmeticError("floating-point stagnation")
        fc = value(c)
        if fc == 0:
            return c, radius, "computed_zero"
        if (fa > 0) != (fc > 0):
            b = c
        else:
            a, fa = c, fc
    raise RuntimeError("iteration limit reached")
\end{verbatim}
理论半区间界适用于精确端点与精确中点；在普通浮点实现中它是位置误差的工作估计，严格机器可验证包围还需向外舍入。相对阈值按当前中点尺度定义；接近零的根主要由绝对容差控制。

\section{复合 Simpson 与标量 RK4}
以下积分代码要求正偶数段数。\texttt{math.fsum} 改善长求和的舍入表现，但它不修复非光滑函数、错误节点或不适用的误差估计。
\begin{verbatim}
def composite_simpson(f, a, b, n):
    if not isinstance(n, int) or n <= 0 or n % 2:
        raise ValueError("n must be a positive even integer")
    if not (math.isfinite(a) and math.isfinite(b) and a < b):
        raise ValueError("finite ordered endpoints required")
    h = (b - a) / n
    if not math.isfinite(h) or h == 0:
        raise ValueError("unusable step size")
    values = [f(a)] + [f(a + j * h) for j in range(1, n)] + [f(b)]
    if not all(math.isfinite(v) for v in values):
        raise ValueError("nonfinite function value")
    odd = math.fsum(values[1:n:2])
    even = math.fsum(values[2:n:2])
    result = h / 3 * math.fsum([values[0], values[-1],
                              4 * odd, 2 * even])
    if not math.isfinite(result):
        raise ArithmeticError("nonfinite quadrature result")
    return result

def rk4_scalar(f, t0, t1, y0, n):
    if not isinstance(n, int) or n <= 0:
        raise ValueError("n must be a positive integer")
    if not all(math.isfinite(v) for v in (t0, t1, y0)):
        raise ValueError("finite data required")
    if t1 <= t0:
        raise ValueError("increasing time interval required")
    h = (t1 - t0) / n
    if not math.isfinite(h) or h == 0:
        raise ValueError("unusable step size")
    y = y0
    for j in range(n):
        t = t0 + j * h
        k1 = f(t, y)
        k2 = f(t + h / 2, y + h * k1 / 2)
        k3 = f(t + h / 2, y + h * k2 / 2)
        k4 = f(t + h, y + h * k3)
        y += h * (k1 + 2 * k2 + 2 * k3 + k4) / 6
        if not all(math.isfinite(v) for v in (k1, k2, k3, k4, y)):
            raise ArithmeticError("nonfinite stage or solution")
    return y
\end{verbatim}
以上函数只返回终点标量值；向量方程需要按分量保存与更新各阶段，而不是把 Python 列表的连接当作向量加法。时间步长由输入固定，没有局部误差自适应，也没有刚性检测，因此使用前应先做第~\ref{sec:na-ode-stability} 节的稳定性判断。

\section{六组综合实验及验收问题}
\subsection{抵消、步长与有效精度}
计算 $g(x)=\sqrt{x^2+1}-x$ 及其有理化形式，取 $x=10^k$、$k=0,1,\ldots,8$；另用中心差分近似 $\exp'(1)=e$，取 $h=10^{-k}$、$k=1,\ldots,16$。用高精度参考值或解析表达式计算误差，记录误差最低点及随后上升或平台。验收时要区分函数评价误差、相消、最终除以步长的放大，不预设每台机器上最低点完全相同。

\subsection{病态系统与不同最小二乘实现}
采用 Hilbert 矩阵 $H_{ij}=1/(i+j-1)$，$1\le i,j\le n$，取 $n=3,5,8,10$，制造精确向量 $x_*=\boldsymbol1$ 并生成 $b=Hx_*$。分别报告求解残差、解误差和条件数估计；形成右端项时也会舍入，应注明是否以更高精度生成参考数据。再对 $m=50$ 个均匀节点上的带扰动数据作多项式拟合，比较单项式基的正规方程、Householder QR 和移位缩放后的正交基。验收关注残差小而系数误差大、基选择改变表示条件等现象，不把所有差异都解释为算法错误。

\subsection{全局插值与分段插值}
对 Runge 函数比较等距节点与 Chebyshev 节点的重心插值，再比较分段线性与自然三次样条。选取多组节点数，在独立且足够密的验证网格上计算最大采样误差；该量是连续一致误差的数值近似，不应声称就是严格最大值。特别观察端点区域、节点值扰动的影响，以及自然边界与原函数端点二阶导数不匹配的现象。

\subsection{求积阶数与非光滑反例}
先对 $e^x$ 在 $[0,1]$ 上用 $N=4,8,16,32,64$ 比较复合梯形、Simpson 和外推，记录误差比。再对 $\sqrt x$ 重复实验，其积分为 $2/3$，但端点导数不满足前述光滑假设。验收要求解释为什么第二组不能机械套用二阶、四阶定理。最后构造一个窄峰函数，说明有限采样的自适应估计可能漏峰。

\subsection{ODE 稳定性与 PDE 网格收敛}
对 $y'=-20y$、$y(0)=1$，比较显式 Euler 在 $h=0.05,0.1,0.15$ 的行为，分别对应衰减、边界振荡和放大。对 $y'=y$ 用 RK4 验证四阶，避免把纯指数测试的正确性当作一般非线性实现的充分验证；另加 $y'=y^2$、$y(0)=1$、$0\le t\le1/2$，精确解为 $1/(1-t)$。Poisson 实验使用制造解 $\sin(\pi x)$，验证二阶空间误差；热方程用 $u(x,t)=e^{-\pi^2t}\sin(\pi x)$，选 $\Delta t=0.4h^2$，检查总误差与稳定性。

\subsection{FFT 约定与随机误差覆盖}
先用手算可核验的四点 DFT 检查频率、符号和归一化，再用补零卷积与直接卷积对比。蒙特卡洛实验对 $\int_0^1x^2\,\mathrm dx$ 采用多组样本量，并在各样本量下重复独立运行，比较经验均方根误差与 $2/\sqrt{45N}$；统计所构造近似置信区间覆盖 $1/3$ 的比例。报告随机数生成器、种子安排、重复次数和区间构造方式，讨论有限重复带来的覆盖率统计波动。

\section{方法选择的统一依据}
若数据已有高噪声，过度细分或增加插值次数可能只是拟合噪声；若矩阵具有正定、稀疏或带状结构，应优先保留该结构；若局部高阶方法被刚性稳定条件迫使使用极小步长，应比较隐式方法的每步求解成本与总步数。需要严格误差证书时，普通浮点差值估计往往不足，应采用具有可证明界的方法或区间算术。

算法选择最终应返回具体数学问题：输入扰动允许多大，输出需要什么误差准则，理论条件是否成立，计算资源限制在哪里。理解这些约束后，再讨论高阶、快速、随机或自适应方法的优势，才能使计算结果既高效又可解释。

\appendix
\chapter{章末习题解答}
\label{app:na-solutions}

\section{第~\ref{chap:na-errors} 章}
\begin{description}
\item[习题~\ref{exr:na-error-rational}] 由半角公式，$1-\cos x=2\sin^2(x/2)$。后式避免从接近 $1$ 的两个数相减，减少已有舍入误差相对小结果的放大；它不消除输入 $x$ 自身的不确定性，也仍受三角函数评价精度影响。
\item[习题~\ref{exr:na-error-norm}] 两列绝对值和分别为 $4,6$，故 $\|A\|_1=6$；两行分别为 $3,7$，故 $\|A\|_\infty=7$；$\|A\|_F=\sqrt{1+4+9+16}=\sqrt{30}$。
\item[习题~\ref{exr:na-error-condition}] 绝对条件数为 $|F'(x)|=1/(2\sqrt x)$，相对条件数为 $|xF'(x)/F(x)|=1/2$。接近零时固定绝对扰动相对于输入可能非常大，而固定相对扰动随输入一起缩小，所以两种尺度下的结果并不矛盾；扰动还须保持输入为正。
\item[习题~\ref{exr:na-error-balance}] 对上界求导得 $2C_1h-C_2u/h^2=0$，故 $h_*=(C_2u/(2C_1))^{1/3}$。二阶导数为正且两端趋于无穷，故它是上界的唯一最小点。上界可能不紧，真实误差还可能含抵消或高阶项，因此不能据此定位真实最小点。
\end{description}

\section{第~\ref{chap:na-roots} 章}
\begin{description}
\item[习题~\ref{exr:na-root-bisect}] 两端函数值为 $-1,5$，存在根；$f'(x)=3x^2-1\ge2$，所以严格递增且根唯一。第 $k$ 个中点从 $k=0$ 起计，误差界为 $2^{-(k+1)}$。$2^{-27}<10^{-8}<2^{-26}$，故至多 $27$ 次中点评价足够。
\item[习题~\ref{exr:na-root-fixed}] 对正数 $x$，算术几何平均不等式给出 $(x+2/x)/2\ge\sqrt2$。故第一步后 $x_k\ge\sqrt2$；此时 $x_{k+1}\le x_k$，序列单调有下界，其正极限满足 $\ell=(\ell+2/\ell)/2$，即 $\ell=\sqrt2$。精确递推为 $x_{k+1}-\sqrt2=(x_k-\sqrt2)^2/(2x_k)$，说明局部二次收敛；已到根时直接结束。
\item[习题~\ref{exr:na-root-multiple}] 对 $x_k\ne1$，普通更新为 $x_{k+1}-1=\tfrac23(x_k-1)$，故线性收敛。修正更新 $x_{k+1}=x_k-3f(x_k)/f'(x_k)=1$，一步到根。若已经在根上，应停止而非求 $0/0$。
\item[习题~\ref{exr:na-root-system}] 设 $F=(x^2+y^2-1,x-y)^{\mathsf T}$，每步求
\[
 \begin{pmatrix}2x_k&2y_k\\1&-1\end{pmatrix}
 \begin{pmatrix}s_1\\s_2\end{pmatrix}
 =-\begin{pmatrix}x_k^2+y_k^2-1\\x_k-y_k\end{pmatrix},
\]
再同时更新两分量。行列式为 $-2(x_k+y_k)$，在 $x_k+y_k=0$ 时奇异。两个真实根为 $(1/\sqrt2,1/\sqrt2)$ 及其相反数，在两根处 Jacobian 均非奇异。
\end{description}

\section{第~\ref{chap:na-direct} 章}
\begin{description}
\item[习题~\ref{exr:na-direct-lu}] $L=\begin{pmatrix}1&0\\2&1\end{pmatrix}$、$U=\begin{pmatrix}2&1\\0&1\end{pmatrix}$。前代得 $y=(1,1)^{\mathsf T}$，回代得 $x=(0,1)^{\mathsf T}$。代回原方程得到 $(1,3)^{\mathsf T}$。
\item[习题~\ref{exr:na-direct-cholesky}] $l_{11}=2$，$l_{21}=1$，$l_{22}=\sqrt{3-1}=\sqrt2$，所以 $L=\begin{pmatrix}2&0\\1&\sqrt2\end{pmatrix}$。直接相乘可恢复 $A$。
\item[习题~\ref{exr:na-direct-condition}] 正交矩阵保二范数，故 $\|Q\|_2=\|Q^{-1}\|_2=1$，条件数为 $1$。给定对角矩阵的最大、最小奇异值为 $1,10^{-6}$，条件数为 $10^6$。
\item[习题~\ref{exr:na-direct-residual}] 用相容范数下的残差界可保证相对前向误差至多 $10^{-5}$。这是最坏方向的上界，不是误差等式；残差可能指向较不敏感的方向，真实误差可远小于它。
\end{description}

\section{第~\ref{chap:na-ls} 章}
\begin{description}
\item[习题~\ref{exr:na-ls-constant}] 目标为 $\sum_i(c-b_i)^2$，求导得 $2mc-2\sum_i b_i=0$，二阶导数 $2m>0$，故 $c=\overline b$。正权情况下目标为 $\sum_iw_i(c-b_i)^2$，唯一极小点为 $c=\sum_iw_ib_i/\sum_iw_i$。
\item[习题~\ref{exr:na-ls-householder}] 令 $P=vv^{\mathsf T}/(v^{\mathsf T}v)$，则 $P^{\mathsf T}=P$、$P^2=P$。因此 $H^{\mathsf T}H=(I-2P)^2=I$，保二范数；又 $Pv=v$，故 $Hv=-v$。
\item[习题~\ref{exr:na-ls-rank}] 所有满足 $x_1+x_2=1$ 的向量都使残差为零，构成全部最小二乘解。写成 $(1/2+t,1/2-t)$，范数平方为 $1/2+2t^2$，故最小范数解为 $(1/2,1/2)^{\mathsf T}$。
\item[习题~\ref{exr:na-ls-projection}] 奇偶性使常数系数为零，$x$ 的系数为 $\langle x^3,x\rangle/\langle x,x\rangle=(2/5)/(2/3)=3/5$。平方误差为 $2/7-(2/5)^2/(2/3)=8/175$。
\end{description}

\section{第~\ref{chap:na-iterative} 章}
\begin{description}
\item[习题~\ref{exr:na-iter-jacobi}] $B_J=\begin{pmatrix}0&-1/2\\-1/2&0\end{pmatrix}$，特征值为 $\pm1/2$，谱半径为 $1/2$；$B_G=\begin{pmatrix}0&-1/2\\0&1/4\end{pmatrix}$，谱半径为 $1/4$。
\item[习题~\ref{exr:na-iter-nonnormal}] 写 $B=\tfrac12I+N$，其中 $N^2=0$，则对 $k\ge1$，$B^k=2^{-k}\begin{pmatrix}1&20k\\0&1\end{pmatrix}$。当 $e_0=(0,1)^{\mathsf T}$ 时，$\|Be_0\|_2=\sqrt{100+1/4}>1$，第一步放大；但 $k2^{-k}\to0$，最终误差仍趋于零。
\item[习题~\ref{exr:na-iter-cg}] $r_0=p_0=(1,1)^{\mathsf T}$，$\alpha_0=2/3$，$x_1=(2/3,2/3)^{\mathsf T}$，$r_1=(1/3,-1/3)^{\mathsf T}$。于是 $\beta_0=1/9$，$p_1=(4/9,-2/9)^{\mathsf T}$，$\alpha_1=3/4$，$x_2=(1,1/2)^{\mathsf T}$，恰为解。
\item[习题~\ref{exr:na-iter-energy}] 由 $r=-A(x-x_*)$，有 $r^{\mathsf T}A^{-1}r=(x-x_*)^{\mathsf T}A(x-x_*)=\|x-x_*\|_A^2$。正定性保证其为非负量并定义范数。
\end{description}

\section{第~\ref{chap:na-eigen} 章}
\begin{description}
\item[习题~\ref{exr:na-eigen-power}] 第 $k$ 次归一化结果为 $(4^k,2^k,1)^{\mathsf T}/\sqrt{16^k+4^k+1}$。第二分量相对第一分量为 $2^{-k}$，第三为 $4^{-k}$，所以主导方向误差量级为 $2^{-k}$。
\item[习题~\ref{exr:na-eigen-rayleigh}] $\mu(v)=\lambda_1\cos^2\theta+\lambda_2\sin^2\theta$，故 $\mu(v)-\lambda_1=(\lambda_2-\lambda_1)\sin^2\theta$。特征向量夹角的一阶小量进入特征值估计时成为二阶。
\item[习题~\ref{exr:na-eigen-shift}] 三个距离为 $1.8,0.2,5.2$，所以在初值第二分量非零的条件下趋向特征值 $3$ 对应的第二坐标轴。其余方向系数相对于该方向依次按距离比 $0.2/1.8$、$0.2/5.2$ 衰减。
\item[习题~\ref{exr:na-eigen-rotation}] 特征多项式为 $\lambda^2+1$，根为 $\pm\mathrm i$。相似变换保持特征值，而实对角矩阵只有实特征值，因此不可能。实 Schur 分解保留此类 $2\times2$ 块。
\end{description}

\section{第~\ref{chap:na-interpolation} 章}
\begin{description}
\item[习题~\ref{exr:na-interp-poly}] Lagrange 形式为 $(x-1)(x-2)/2-2x(x-2)+5x(x-1)/2$。差商依次为 $1,1,1$，Newton 形式为 $1+x+x(x-1)$。两者都化为 $x^2+1$，故 $p(3/2)=13/4$。
\item[习题~\ref{exr:na-interp-error}] 最大二阶导数为 $e$，$h=1/N$，故误差至多 $e/(8N^2)$。只要 $N\ge\lceil\sqrt{e/(8\times10^{-4})}\rceil=59$ 即足够。这是充分条件，未声称是最小可行段数。
\item[习题~\ref{exr:na-interp-spline}] 自然边界给出 $M_0=M_2=0$，唯一内部方程为 $4M_1=6(-1-1)=-12$，故 $M_1=-3$。第一段 $0\le x\le1$ 上 $S(x)=-x^3/2+3x/2$，第二段由关于 $x=1$ 对称得到。
\item[习题~\ref{exr:na-interp-cheb}] $T_2=2x^2-1$，$T_3=4x^3-3x$，零点为 $0,\pm\sqrt3/2$。首一化后为 $x^3-3x/4=T_3/4$，最大绝对值为 $1/4$，与 $n+1=3$ 时的 $2^{-n}$ 一致。
\end{description}

\section{第~\ref{chap:na-quadrature} 章}
\begin{description}
\item[习题~\ref{exr:na-quad-difference}] 分别展开 $f(x+h)$ 与 $f(x-h)$ 至三阶并保留四阶 Lagrange 余项。相加减去 $2f(x)$ 后，奇次项消失，除以 $h^2$ 得 $f''(x)+h^2[f^{(4)}(\xi_+)+f^{(4)}(\xi_-)]/24$，故误差不超过 $h^2\|f^{(4)}\|_\infty/12$。
\item[习题~\ref{exr:na-quad-simpson}] 对 $x^3$，$S=(0+4/8+1)/6=1/4=I$。对 $x^4$，$S=(0+4/16+1)/6=5/24$，而 $I=1/5$，所以 $I-S=-1/120$。公式 $-(b-a)^5f^{(4)}/2880=-24/2880$ 给出同一值。
\item[习题~\ref{exr:na-quad-gauss}] 设节点为 $\pm c$、共同权重为 $w$。常数精确要求 $2w=2$，故 $w=1$；二次项要求 $2c^2=2/3$，故 $c=1/\sqrt3$。一、三次项因对称性自动精确。
\item[习题~\ref{exr:na-quad-richardson}] $(4A(h/2)-A(h))/3=A-dh^4/4+O(h^6)$。二阶项被消去，四阶系数由 $(4/16-1)/3=-1/4$ 得到。
\end{description}

\section{第~\ref{chap:na-ode} 章}
\begin{description}
\item[习题~\ref{exr:na-ode-euler}] 显式值为 $(1+1/N)^N<e$。隐式值为 $(1-1/N)^{-N}>e$，$N>1$。分别对数并使用 $\log(1+x)<x$、$-\log(1-x)>x$ 证明。两者都趋于 $e$，但来自不同方向。
\item[习题~\ref{exr:na-ode-stability}] 要求 $|1-20h|\le1$，正步长范围为 $0<h\le0.1$。在 $h=0.1$ 时放大因子为 $-1$，数值解不衰减而交替变号；若希望非负且不振荡，还需 $0<h\le0.05$，在右端第一步即变为零。
\item[习题~\ref{exr:na-ode-rk}] 令 $z=h\lambda$，逐个代入四阶段可得 $R(z)=1+z+z^2/2+z^3/6+z^4/24$。它与 $e^z$ 匹配到四次项，但只证明在线性测试方程上的展开，不能单独作为一般非线性 RK4 四阶性的证明。
\item[习题~\ref{exr:na-ode-root}] 取 $\alpha=(2,-3,1)$、$\beta=(-1,0,0)$，有 $\sum\alpha_j=0$、$\sum j\alpha_j=-1=\sum\beta_j$，故一致。$\rho(\zeta)=\zeta^2-3\zeta+2=(\zeta-1)(\zeta-2)$ 存在模大于 $1$ 的根，因此非零稳定。即使对 $y'=0$，起步误差也可激发 $2^n$ 模态；例如 $y_0=0,y_1=h$ 时 $y_n=h(2^n-1)$，固定正时间、$n\approx t/h$ 时发散。
\end{description}

\section{第~\ref{chap:na-pde} 章}
\begin{description}
\item[习题~\ref{exr:na-pde-exact}] 精确解为 $u(x)=x(1-x)$。二次多项式的二阶中心差分等于精确二阶导数，故采样值满足全部离散方程和边界条件；由离散系统唯一性，数值解在节点上等于采样值。浮点实现仍可有代数舍入误差。
\item[习题~\ref{exr:na-pde-spectrum}] 用 $\sin(\theta+\varphi)+\sin(\theta-\varphi)=2\sin\theta\cos\varphi$，得到 $\lambda_k=2h^{-2}(1-\cos(k\pi/N))=4h^{-2}\sin^2(k\pi/(2N))$。这些特征值均正；最大值与 $h^{-2}$ 同阶，最小值趋于 $\pi^2$。
\item[习题~\ref{exr:na-pde-stability}] 该模态的相邻值与自身相反，放大因子为 $1-4r=-1.4$，故第 $n$ 步为 $(-1.4)^n(-1)^i$，振荡且绝对值指数增长。
\item[习题~\ref{exr:na-pde-mesh}] 取 $\Delta t=ch^2$、固定 $c>0$ 可使 $O(\Delta t+h^2)=O(h^2)$。还必须满足 $\nu c\le1/2$ 的稳定充分条件，并检查初值、边界及解的光滑性条件。仅平衡误差而忽略稳定限制是不完整的设计。
\end{description}

\section{第~\ref{chap:na-fourier} 章}
\begin{description}
\item[习题~\ref{exr:na-fft-transform}] 只有 $j=0,2$ 非零，故 $\widehat v_k=1-e^{-\pi\mathrm i k}=1-(-1)^k$，变换为 $(0,2,0,2)$。逆变换得 $v_j=(e^{\pi\mathrm i j/2}+e^{3\pi\mathrm i j/2})/2$，在四个整数节点依次为 $1,0,-1,0$。
\item[习题~\ref{exr:na-fft-alias}] $7\equiv-3\pmod{10}$，故样本无法区分频率 $7$ 与 $-3$。这不意味着连续函数本身相同。
\item[习题~\ref{exr:na-fft-convolution}] 线性卷积为 $(3,10,8)$；长度为 $2$ 的周期卷积把下标 $2$ 绕回下标 $0$，得到 $(11,10)$。至少补零到长度 $3$ 才能避免绕回，取 FFT 友好的长度 $4$ 也可。
\item[习题~\ref{exr:na-fft-periodic}] 由离散正交性，若 $N\nmid m$，节点平均为零，积分近似为 $0$；若 $N\mid m$，各节点余弦均为 $1$，结果为 $2\pi$。真实积分在 $m\ge1$ 时均为零，后一情形体现完全混叠。
\end{description}

\section{第~\ref{chap:na-montecarlo} 章}
\begin{description}
\item[习题~\ref{exr:na-mc-variance}] 估计器为样本平均 $\overline X$，期望 $1/2$，方差 $1/(12N)$，均方根误差 $1/\sqrt{12N}$。这里方差来自独立均匀样本，不能在相关样本下直接沿用。
\item[习题~\ref{exr:na-mc-antithetic}] 每一对的平均为 $[U+(1-U)]/2=1/2$，因此方差为零。这个特殊结构并不说明任意函数的对偶变量法都零方差。
\item[习题~\ref{exr:na-mc-control}] $\operatorname{Cov}(U^2,U)=1/4-(1/3)(1/2)=1/12$，$\operatorname{Var}(U)=1/12$，故 $c_*=1$。修正量为 $U^2-U+1/2$，方差为 $4/45-(1/12)^2/(1/12)=1/180$，比原方差 $4/45$ 小为原来的 $1/16$。
\item[习题~\ref{exr:na-mc-importance}] 单样本权重为 $Y=1/(2X)$，其期望 $\int_0^1[1/(2x)]2x\,\mathrm dx=1$，但二阶矩为 $\int_0^1[1/(4x^2)]2x\,\mathrm dx=\tfrac12\int_0^1x^{-1}\,\mathrm dx=\infty$。因此无偏性并不蕴含有限方差，常规标准误差与有限方差中心极限定理的公式不适用。
\end{description}

\chapter{公式使用前的条件核对}
\label{app:na-conditions}

\section{存在性、误差与计算保证的区别}
求根中的连续异号保证根存在，导数非零保证局部 Newton 行为，两者承担不同作用；最小二乘解存在不意味着唯一，列满秩才保证参数唯一；正定性保证 Cholesky 和 CG 的核心结构，却不意味着任意步长的时间积分稳定。检查条件时，应将所使用定理的对象与实际输入逐项对应。

误差界中的常数常依赖导数上界、区间长度或条件数。一个形式上为 $O(h^4)$ 的方法，在四阶导数很大时可能不如常数较小的二阶方法；如果函数不满足四阶光滑条件，则该阶数甚至没有适用依据。局部缺陷、全局误差、代数残差和随机标准误差分别描述不同量，不应在图表或报告中只写一个未定义的“误差”。

\section{正文未完整展开的理论}
为保持本科主线，正文对奇异值分解给出了基于谱定理的证明要点，对 CG 给出投影与多项式误差界的逻辑，对割线法阶数、一般 Runge--Kutta 高阶条件、线性多步法等价定理和一般 QR 迭代仅说明结论或证明范围。进一步学习应专门补足这些证明及浮点误差分析。正文中已展开的基础证明包括压缩迭代、Newton 局部误差、线性系统扰动、Cholesky、最小二乘正规条件、插值余项、Gauss 精确性、Euler 全局界和一维差分稳定收敛。

\begin{thebibliography}{9}
\bibitem{yu2020}
喻文健. 数值分析与算法：第3版. 清华大学出版社，2020.
ISBN 9787302544616.
本书核验范围为出版社图书介绍及公开目录：
\url{https://www.tup.tsinghua.edu.cn/bookscenter/book_08630401.html}.

\bibitem{pku-course}
北京大学数学科学学院. 本科常设课程：数值分析，课程号 00130560.
公开课程大纲：\url{https://math.pku.edu.cn/bks/sykc/148732.htm}.

\bibitem{mit-syllabus}
Laurent Demanet. 18.330 Introduction to Numerical Analysis, Spring 2012: Syllabus.
Massachusetts Institute of Technology, MIT OpenCourseWare.
\url{https://ocw.mit.edu/courses/18-330-introduction-to-numerical-analysis-spring-2012/pages/syllabus/}.

\bibitem{mit-notes}
Laurent Demanet. 18.330 Introduction to Numerical Analysis, Spring 2012: Lecture Notes.
MIT OpenCourseWare. 此处引用课程讲义主题索引，用于课程范围比较与延伸阅读：
\url{https://ocw.mit.edu/courses/18-330-introduction-to-numerical-analysis-spring-2012/pages/lecture-notes/}.
\end{thebibliography}

% 公开课程与出版社元数据核验日期：2026-10-11。
\end{document}
